\documentclass[11pt]{article}

\usepackage[margin=1in]{geometry}
\usepackage{amsmath,amssymb,amsthm,mathtools}
\usepackage{booktabs,tabularx}
\usepackage[authoryear]{natbib}
\usepackage[colorlinks=true,allcolors=blue]{hyperref}
\usepackage{microtype}

\newtheorem{theorem}{Theorem}[section]
\newtheorem{proposition}[theorem]{Proposition}
\newtheorem{corollary}[theorem]{Corollary}
\newtheorem{lemma}[theorem]{Lemma}
\theoremstyle{definition}
\newtheorem{definition}[theorem]{Definition}
\theoremstyle{remark}
\newtheorem{remark}[theorem]{Remark}

\newcommand{\Acal}{\mathcal{A}}
\newcommand{\Ccal}{\mathcal{C}}
\newcommand{\Ocal}{\mathcal{O}}
\newcommand{\Pcal}{\mathcal{P}}
\newcommand{\Fcal}{\mathcal{F}}
\newcommand{\Gcal}{\mathcal{G}}
\newcommand{\Qcal}{\mathcal{Q}}
\newcommand{\Rcal}{\mathcal{R}}
\newcommand{\E}{\mathbb{E}}
\newcommand{\Prob}{\mathbb{P}}
\newcommand{\KL}{\mathrm{KL}}
\newcommand{\I}{\mathrm{I}}
\newcommand{\Hh}{\mathrm{H}}
\newcommand{\argmin}{\mathop{\mathrm{arg\,min}}}
\newcommand{\argmax}{\mathop{\mathrm{arg\,max}}}
\newcommand{\givena}{\,\|\,}
\newcommand{\ellp}{\ell}
\newcommand{\wt}{\widetilde}
\newcommand{\barxi}{\bar{\xi}}
\newcommand{\one}{\mathbf{1}}
\newcommand{\problembox}[1]{%
\begin{center}
\fbox{\parbox{0.93\textwidth}{#1}}
\end{center}}

\title{Sufficient Conditions for Semantic Convergence in Interactive Learning}
\author{}
\date{}

\begin{document}

\maketitle

\begin{abstract}
We seek sufficient conditions under which an interactive learner's posterior concentrates on the correct \emph{interventional semantic class} $[p^\star]$ of an unknown computable environment --- the programs that induce the same action-conditioned law as $p^\star$ under every admissible intervention, and the strongest target any history-based learner can reach. A natural means is the \emph{algorithmic free energy}: the Gibbs-form variational functional whose reference measure is the Kolmogorov-complexity-weighted universal interactive prior of \citet{Solomonoff1964a,Solomonoff1964b,Kolmogorov1965,Hutter2005}, rather than the stipulated parametric prior of active-inference free energy. On this functional, two unifications are exact: minimizing it is Bayesian updating, with the KL-based divergence forced within the natural convex class; and with zero action cost and uniform preferences, the resulting expected free energy is the information-gain criterion shared by Bayesian experimental design, compression-based curiosity, and active-inference curiosity. The trio --- universal prior, Bayes/Gibbs belief, expected-free-energy-minimizing action --- is the \emph{algorithmic free energy principle}. Perhaps unintuitively, even the same Bayes/Gibbs/expected-free-energy architecture need not suffice: a near-miss theorem for a substituted finite-support prior shows that the posterior on $[p^\star]$ need not concentrate at any horizon, and because the action rule there is exactly the unifying generic information-gain criterion, the failure is a failure of generic information-seeking itself, not of a peculiar implementation choice. Concentration to the correct belief requires a class-against-complement action primitive: under the universal prior, Bayes/Gibbs belief, a semantic separation condition, and the class-against-complement rule, the posterior mass on $[p^\star]$ converges almost surely to one.
\end{abstract}

\section{Introduction}
\label{sec:intro}

Prediction admits many successes that explanation does not. An interactive learner may reproduce a realized transcript, forecast the next observation along the policy it happens to follow, and reduce uncertainty about the latent program in general, and still be wrong about what the environment would do under an intervention it never tried. In an interactive problem, that residual ambiguity is the question worth asking: whether the learner's internal inference converges to the environment law that is correct under every admissible intervention.

Four targets organize the answer for computable interactive environments. The \emph{history-compatible set} $R_{h_t}(p^\star)$ consists of programs that reproduce the realized transcript; the \emph{policy-predictable set} $R_\pi(p^\star)$, of those that induce the same law as $p^\star$ under the policy the learner actually follows; the \emph{interventional semantic class} $[p^\star]$, of those that induce the same action-conditioned law under every admissible intervention; and the innermost target, the singleton $\{p^\star\}$, is the truth. These form a nested hierarchy,
\[
\{p^\star\}
\;\subseteq\;
[p^\star]
\;\subseteq\;
R_\pi(p^\star)
\;\subseteq\;
R_{h_t}(p^\star)
\;\subseteq\;
\Pcal,
\]
and all three inclusions between the truth and the full compatibility set are strict in general (Theorem~\ref{thm:strict-hierarchy}). An observable-quotient theorem (Theorem~\ref{thm:factor-through-quotient}) makes $[p^\star]$ the theorem-level ceiling: any target a history-based learner can reliably recover is constant on $[p^\star]$, so concentration on $\{p^\star\}$ itself is impossible whenever a semantic class contains more than one program. The four sets are the indistinguishability classes of four \emph{observers} at increasing observational power (Definition~\ref{def:observer}), from reading only the realized transcript up to reading the source code; the paper's theorems turn on whether a learning rule \emph{promotes} the realized observer toward the behavioral one.

Three existing traditions supply natural ingredients for a learner aimed at that ceiling. Solomonoff induction and AIXI supply a universal prior over computable interactive environments, weighted by Kolmogorov complexity \citep{Solomonoff1964a,Solomonoff1964b,Kolmogorov1965,Hutter2005}, but no variational functional. The free energy principle and active inference supply the Gibbs variational form and a risk-plus-epistemic-value action rule \citep{Friston2010,FristonEtAl2017Curiosity,ParrFriston2019}, but with a stipulated parametric prior whose justification has been debated on philosophical grounds \citep{ColombWright2021}. Bayesian experimental design, information-theoretic bounded rationality, and compression-based curiosity supply closely related information-seeking action criteria, again without a universal-prior reference measure \citep{Lindley1956,ChalonerVerdinelli1995,MacKay1992,OrtegaBraun2013,Schmidhuber1991,Schmidhuber2010}. No existing framework couples these three ingredients into a single variational object. We formulate the \emph{algorithmic free energy} --- the Gibbs-form variational functional whose reference measure is the Kolmogorov-complexity-weighted universal interactive prior rather than a stipulated parametric prior (Definition~\ref{def:afe}) --- and on it the three traditions unify on both sides of the inference problem. The Bayes/Gibbs posterior is its unique minimizer (Lemma~\ref{lem:variational}), and the KL-based divergence is forced within the natural convex class (Lemma~\ref{lem:kl-necessity}); minimizing the algorithmic free energy is not an approximation of Bayesian updating but is Bayesian updating exactly. Under zero action cost and uniform preferences, the resulting expected free energy coincides with the information-gain criterion of Bayesian experimental design, compression-based curiosity, and active-inference curiosity (Corollary~\ref{cor:efe-specialization}).

The unifications would appear to close the problem. They do not. The same Bayes/Gibbs/expected-free-energy architecture can fail to concentrate its posterior on $[p^\star]$ (Theorem~\ref{thm:afe-near-miss}): for every horizon $T$ and every nontrivial prior mass $\alpha_0\in(0,1/2)$ there is a problem instance, using a substituted finite-support prior, on which the posterior on $[p^\star]$ is frozen at $\alpha_0$ for all $t\le T$. The theorem extends the AIXI near-miss tradition \citep{LeikeHutter2015,Leike2016,LeikeEtAl2016} and locates the failure specifically at the level of the interventional semantic class. Because expected-free-energy action is exactly the unifying generic information-gain criterion, the failure is a failure of generic information-seeking itself, not of any peculiar implementation choice. In observer language, the near-miss construction makes the AFE principle's realized policy induce an observer strictly coarser than the behavioral observer, and the frozen-posterior conclusion is the Bayesian shadow of that strict refinement (Theorem~\ref{thm:observer-promotion-failure}).

\problembox{\textbf{Problem.} Give sufficient conditions on the learner's prior, belief rule, and action rule under which the posterior mass on $[p^\star]$ converges almost surely to one.}

The answer is a class-against-complement action primitive: instead of reducing uncertainty about the latent program in general, the learner increases the posterior odds of a semantic class against the aggregate posterior mass on every other class (Section~\ref{sec:semantic-action}). Paired with the universal prior, Bayes/Gibbs belief, and an explicit semantic separation condition, this primitive suffices for almost-sure convergence (Theorem~\ref{thm:semantic-convergence}, Section~\ref{sec:main}).

Section~\ref{sec:setup} fixes the interactive setting. Section~\ref{sec:hierarchy} constructs the four-set hierarchy, proves the observable-quotient theorem, and witnesses the three strict inclusions. Section~\ref{sec:belief} proves the belief-side unification. Section~\ref{sec:near-miss} proves the action-side unification and the near-miss. Section~\ref{sec:semantic-action} introduces the class-against-complement action primitive. Section~\ref{sec:main} proves the main sufficient-conditions theorem. Sections~\ref{sec:implications},~\ref{sec:discussion}, and~\ref{sec:conclusion} draw consequences.

\section{Formal Setup}
\label{sec:setup}

Let $\Acal$ be a countable action alphabet and $\Ocal$ a countable observation alphabet. Time is discrete. At step $t$ the learner emits an action $a_t\in\Acal$ and receives an observation $o_t\in\Ocal$. Histories are
\[
h_t=(a_1,o_1,\dots,a_t,o_t),
\qquad
h_0=\varepsilon.
\]
A policy is a sequence of kernels $\pi_t(a_t\mid h_{t-1})$. Actions are interventions and observations are evidence. Under an environment $p$ and a policy $\pi$, write $(A_t,O_t)_{t\ge 1}$ for the resulting action-observation process and
\[
H_t:=(A_1,O_1,\dots,A_t,O_t)
\]
for the random history at time $t$.

Fix a universal prefix machine $U$. Let $\Pcal$ be the countable set of programs in the prefix-free domain of $U$ that induce interactive environment models. Each $p\in\Pcal$ defines an action-conditioned chronological law
\[
\mu_p(o_{1:t}\givena a_{1:t})
=
\prod_{i=1}^{t}\mu_p(o_i\mid h_{i-1},a_i),
\qquad
\sum_{o\in\Ocal}\mu_p(o\mid h_{t-1},a_t)=1.
\]
Throughout, kernel equality $\mu_p=\mu_{p'}$ denotes agreement at reachable pairs: $\mu_p(\cdot\mid h,a)=\mu_{p'}(\cdot\mid h,a)$ whenever $\mu_p(o_{1:|h|}\givena a_{1:|h|})>0$ or $\mu_{p'}(o_{1:|h|}\givena a_{1:|h|})>0$; jointly unreachable pairs have no observable consequence under any policy.

For a prior $w$ on $\Pcal$ and a history $h_t$ with positive $w$-induced marginal likelihood, the Bayes posterior over programs is
\[
q_t(p\mid h_t)
:=
\frac{w(p)\mu_p(o_{1:t}\givena a_{1:t})}
{\sum_{r\in\Pcal} w(r)\mu_r(o_{1:t}\givena a_{1:t})}.
\]
For any subset $\Rcal\subseteq\Pcal$ write
\[
q_t(\Rcal\mid h_t):=\sum_{p\in\Rcal} q_t(p\mid h_t).
\]
For a continuation $(a,o)$ after $h_t$ with positive predictive normalizing constant,
\[
q_{t+1}(p\mid h_t,a,o)
:=
\frac{q_t(p\mid h_t)\mu_p(o\mid h_t,a)}
{\sum_{r\in\Pcal} q_t(r\mid h_t)\mu_r(o\mid h_t,a)}.
\]

\section{The Knowability Hierarchy}
\label{sec:hierarchy}

This section replaces a list of informal learning objectives with four nested sets of programs and proves that each inclusion is strict. The strictness results are the structural spine of the paper.

\subsection{Four nested sets}

\begin{definition}[History-compatible set]
\label{def:history-compat}
For a history $h_t$, the \emph{history-compatible set} is
\[
R_{h_t}(p^\star)
:=
\bigl\{p\in\Pcal\,:\,\mu_p(o_{1:t}\givena a_{1:t})=\mu_{p^\star}(o_{1:t}\givena a_{1:t})\bigr\}.
\]
\end{definition}

\begin{definition}[Policy-predictable set]
\label{def:policy-pred}
For a policy $\pi$, the \emph{policy-predictable set} is
\[
R_\pi(p^\star)
:=
\bigl\{p\in\Pcal\,:\,\Prob_{p,\pi}(H_T\in\cdot)=\Prob_{p^\star,\pi}(H_T\in\cdot)\text{ for every }T\bigr\}.
\]
\end{definition}

\begin{definition}[Interventional semantic class]
\label{def:int-sem-class}
The \emph{interventional semantic class} of $p^\star$ is
\[
[p^\star]
:=
\bigl\{p\in\Pcal\,:\,\mu_p=\mu_{p^\star}\bigr\}.
\]
\end{definition}

Write $p\equiv p'$ for semantic equivalence and $\Ccal:=\Pcal/\!\equiv$ for the set of classes. Two programs lie in the same class exactly when they agree on the entire action-conditioned environment law.

\begin{definition}[Observer]
\label{def:observer}
An \emph{observer} is a pair $\omega=(\Omega_\omega,V_\omega)$ of a measurable space $\Omega_\omega$ and a measurable view map $V_\omega:\Pcal\to\Omega_\omega$. Observer $\omega$ is \emph{refined by} $\omega'$, written $\omega\le\omega'$, if there is a measurable factorization $\phi:\Omega_{\omega'}\to\Omega_\omega$ with $V_\omega=\phi\circ V_{\omega'}$. The \emph{$\omega$-indistinguishability class} of $p^\star$ is the fiber $V_\omega^{-1}(V_\omega(p^\star))$: the programs an adversary could substitute for $p^\star$ without the $\omega$-observer detecting the swap.
\end{definition}

Four canonical observers structure the hierarchy: the \emph{syntactic observer} $\omega_{\mathrm{syn}}$ with $V_{\omega_{\mathrm{syn}}}(p):=p$, reading the source code; the \emph{behavioral observer} $\omega_{\mathrm{behav}}$ with $V_{\omega_{\mathrm{behav}}}(p):=\mu_p$, running every policy and seeing every induced law; the \emph{policy observer} $\omega_\pi$ with $V_{\omega_\pi}(p):=\Prob_{p,\pi}$, seeing only the law under a fixed policy; and the \emph{history observer} $\omega_{h_t}$ with $V_{\omega_{h_t}}(p):=\mu_p(o_{1:t}\givena a_{1:t})$, reading only the realized transcript. Their indistinguishability classes of $p^\star$ are exactly $\{p^\star\}$, $[p^\star]$, $R_\pi(p^\star)$, and $R_{h_t}(p^\star)$.

\begin{remark}
\label{rem:four-observers}
The hierarchy is the partition-refinement chain of four observers at increasing observational power: as the observer is promoted, the indistinguishability class shrinks. The nesting, strict-separation, and observable-quotient results below are the three consequences of this single reading. The near-miss pair (Theorems~\ref{thm:afe-near-miss} and~\ref{thm:observer-promotion-failure}) and the main theorem (Theorem~\ref{thm:semantic-convergence}) turn on whether a learning rule \emph{promotes} the realized policy observer $\omega_\pi$ to the behavioral observer $\omega_{\mathrm{behav}}$: the former show that a natural rule does not, the latter gives a sufficient condition under which a stricter rule does.
\end{remark}

\begin{lemma}[Nesting]
\label{lem:nesting}
For any $p^\star\in\Pcal$, any policy $\pi$, and any history $h_t$ with $\Prob_{p^\star,\pi}(H_t=h_t)>0$,
\[
\{p^\star\}\subseteq [p^\star]\subseteq R_\pi(p^\star)\subseteq R_{h_t}(p^\star)\subseteq \Pcal.
\]
\end{lemma}

\begin{proof}
A swap the stronger observer cannot detect, the weaker observer cannot detect. If $p\in[p^\star]$ then $\mu_p=\mu_{p^\star}$, so the induced $\pi$-law matches $p^\star$'s, giving $[p^\star]\subseteq R_\pi(p^\star)$. If $p\in R_\pi(p^\star)$ then $\Prob_{p,\pi}(H_t=h_t)=\Prob_{p^\star,\pi}(H_t=h_t)$; combining with the factorization
\[
\Prob_{p,\pi}(H_t=h_t)=\mu_p(o_{1:t}\givena a_{1:t})\prod_{i=1}^t \pi_i(a_i\mid h_{i-1}),
\]
and dividing through by the policy factor (positive by assumption), one gets $\mu_p(o_{1:t}\givena a_{1:t})=\mu_{p^\star}(o_{1:t}\givena a_{1:t})$, hence $p\in R_{h_t}(p^\star)$. The outer inclusion is immediate.
\end{proof}

\begin{proposition}[Refinement chain]
\label{prop:refinement-chain}
Under the hypotheses of Lemma~\ref{lem:nesting},
\[
\omega_{h_t}\le\omega_\pi\le\omega_{\mathrm{behav}}\le\omega_{\mathrm{syn}}.
\]
\end{proposition}

\begin{proof}
The factorization $p\mapsto\mu_p$ witnesses $\omega_{\mathrm{behav}}\le\omega_{\mathrm{syn}}$. The history-factorization identity of Lemma~\ref{lem:nesting} exhibits $\Prob_{p,\pi}$ as a fixed measurable function of $\mu_p$, giving $\omega_\pi\le\omega_{\mathrm{behav}}$. Positivity of the policy factor along $h_t$ lets one recover $\mu_p(o_{1:t}\givena a_{1:t})=\Prob_{p,\pi}(H_t=h_t)/\prod_i\pi_i(a_i\mid h_{i-1})$, giving $\omega_{h_t}\le\omega_\pi$.
\end{proof}

\subsection{Observable quotient: why $[p^\star]$ is the target}

\begin{lemma}[Interventional semantic classes are the observable quotient]
\label{lem:observable-quotient}
If $p\equiv p'$, then for every policy $\pi$ and every horizon $T$, the induced law of $H_T$ under $(p,\pi)$ equals the induced law of $H_T$ under $(p',\pi)$.
\end{lemma}

\begin{proof}
Fix $h_T$. If $\Prob_{p,\pi}(H_T=h_T)>0$ or $\Prob_{p',\pi}(H_T=h_T)>0$, then every prefix $(h_{t-1},a_t)$ visited along $h_T$ is reachable under $p$ or $p'$, so semantic equivalence gives $\mu_p(o_t\mid h_{t-1},a_t)=\mu_{p'}(o_t\mid h_{t-1},a_t)$ for each $t$. Using the factorization
\[
\Prob_{p,\pi}(H_T=h_T)
=
\prod_{t=1}^{T}\pi_t(a_t\mid h_{t-1})\,\mu_p(o_t\mid h_{t-1},a_t),
\]
and the analogous identity for $p'$, the two probabilities are equal. If both probabilities are zero, equality is trivial. Hence the laws coincide.
\end{proof}

\begin{definition}[History-recoverable target]
\label{def:history-recoverable}
Let $(\mathcal{T},d)$ be a metric space and let $\tau:\Pcal\to\mathcal{T}$ be a target map. The target $\tau$ is \emph{history-recoverable} if there exist a history-based policy $\pi$ and measurable estimators $\hat\tau_t:(\Acal\times\Ocal)^t\to\mathcal{T}$ such that for every $p\in\Pcal$,
\[
\hat\tau_t(H_t)\longrightarrow \tau(p)
\qquad
\text{in probability under the law generated by } (p,\pi).
\]
\end{definition}

\begin{theorem}[Every history-recoverable target factors through the semantic quotient]
\label{thm:factor-through-quotient}
Let $(\mathcal{T},d)$ be a metric space and let $\tau:\Pcal\to\mathcal{T}$ be history-recoverable. Then $p\equiv p'\Rightarrow\tau(p)=\tau(p')$. Equivalently, there exists a unique map $\bar\tau:\Ccal\to\mathcal{T}$ such that $\tau=\bar\tau\circ[\cdot]$.
\end{theorem}

\begin{proof}
Fix $p\equiv p'$. Let $\pi$ and $\hat\tau_t$ witness history-recoverability, and set $Y_t:=\hat\tau_t(H_t)$. By Lemma~\ref{lem:observable-quotient}, the law of $H_t$ under $(p,\pi)$ equals the law of $H_t$ under $(p',\pi)$ for every $t$. Since $Y_t$ is a measurable function of $H_t$, the law of $Y_t$ is the same under $(p,\pi)$ and $(p',\pi)$. Definition~\ref{def:history-recoverable} gives $Y_t\to\tau(p)$ in probability under $(p,\pi)$ and $Y_t\to\tau(p')$ in probability under $(p',\pi)$. Weak limits in a metric space are unique, so $\delta_{\tau(p)}=\delta_{\tau(p')}$ and $\tau(p)=\tau(p')$. Define $\bar\tau([p]):=\tau(p)$.
\end{proof}

\begin{remark}
Theorem~\ref{thm:factor-through-quotient} is the universality of the semantic quotient expressed in the adversarial reading of Remark~\ref{rem:four-observers}: no history-based estimator can resolve distinctions finer than the full-policy observer detects. Concentration on a target finer than $[p^\star]$ --- in particular, on $\{p^\star\}$ --- asks the observer to see through what the adversary can hide by substituting a semantically equivalent twin.
\end{remark}

\subsection{Strict separation of the hierarchy}

\begin{lemma}[{Fit gap: $R_\pi(p^\star)\subsetneq R_{h_t}(p^\star)$}]
\label{lem:fit-gap}
Assume $|\Ocal|\ge 2$, fix any policy $\pi$, and fix any history $h_t$ with $\Prob_{p^\star,\pi}(H_t=h_t)>0$ and any next action $a$ with $\pi_{t+1}(a\mid h_t)>0$ such that $\mu_{p^\star}(\cdot\mid h_t,a)$ is not a point mass at a single observation in $\Ocal$. Then there exists $p\in R_{h_t}(p^\star)\setminus R_\pi(p^\star)$.
\end{lemma}

\begin{proof}
Let $x\in\Ocal$ be any observation with $\mu_{p^\star}(x\mid h_t,a)<1$. Construct a computable environment $p$ that agrees with $p^\star$ on $\mu(o_{1:t}\givena a_{1:t})$ for the fixed interaction prefix $h_t$ (achievable by copying $p^\star$'s history factorization on that prefix) and that emits $x$ deterministically at $(h_t,a)$, extending arbitrarily elsewhere. Then $\mu_p(o_{1:t}\givena a_{1:t})=\mu_{p^\star}(o_{1:t}\givena a_{1:t})$, so $p\in R_{h_t}(p^\star)$; and $\mu_p(\cdot\mid h_t,a)\neq \mu_{p^\star}(\cdot\mid h_t,a)$, so under $\pi$ the induced law of $H_{t+1}$ on the event $\{H_t=h_t,A_{t+1}=a\}$ differs between $p$ and $p^\star$, so $p\notin R_\pi(p^\star)$.
\end{proof}

\begin{theorem}[{Policy gap: $[p^\star]\subsetneq R_\pi(p^\star)$}]
\label{thm:policy-gap}
Assume $|\Acal|\ge 2$. Let $\pi$ be any policy. Suppose there exist a history $h_{t_0}$ with $\Prob_{p^\star,\pi}(H_{t_0}=h_{t_0})>0$, an action $b\in\Acal$ with $\pi_{t_0+1}(b\mid h_{t_0})=0$, and an observation $y\in\Ocal$ with $\mu_{p^\star}(y\mid h_{t_0},b)<1$. Then there exists $p'\in R_\pi(p^\star)\setminus [p^\star]$.
\end{theorem}

\begin{proof}
Construct a computable environment $p'$ that agrees with $p^\star$ at every conditioning pair $(h,a)\neq(h_{t_0},b)$ and that emits $y$ deterministically at $(h_{t_0},b)$. Such a $p'$ exists: if $p^\star$ has a program in $\Pcal$, adding a fixed conditional override at a single pair yields another program in $\Pcal$.

Because $\pi_{t_0+1}(b\mid h_{t_0})=0$, any history that visits the pair $(h_{t_0},b)$ has probability zero under every program in $\Pcal$ paired with $\pi$. For a history $h_T$ that does not visit $(h_{t_0},b)$, every conditioning pair along $h_T$ is unchanged, so
\[
\Prob_{p',\pi}(H_T=h_T)
=
\prod_{t=1}^T \pi_t(a_t\mid h_{t-1})\mu_{p'}(o_t\mid h_{t-1},a_t)
=
\prod_{t=1}^T \pi_t(a_t\mid h_{t-1})\mu_{p^\star}(o_t\mid h_{t-1},a_t)
=
\Prob_{p^\star,\pi}(H_T=h_T),
\]
because every factor $\mu_{p'}(o_t\mid h_{t-1},a_t)$ along such a history agrees with $\mu_{p^\star}(o_t\mid h_{t-1},a_t)$. Thus the two laws agree on every history, and $p'\in R_\pi(p^\star)$.

On the other hand, $\mu_{p'}(\cdot\mid h_{t_0},b)=\delta_y\neq \mu_{p^\star}(\cdot\mid h_{t_0},b)$, so $\mu_{p'}\neq\mu_{p^\star}$ and $p'\notin [p^\star]$.
\end{proof}

\begin{remark}
Theorem~\ref{thm:policy-gap} is the central separation in the hierarchy. The constructed $p'$ is the adversary's move: it agrees with $p^\star$ on every pair the $\pi$-fixed observer sees and disagrees only at the pair that observer never queries. Closing this gap is therefore a matter of policy, not of belief; no deterministic or support-limited policy that leaves a semantically informative reachable pair unqueried can close it.
\end{remark}

\begin{lemma}[{Syntactic gap: $\{p^\star\}\subsetneq [p^\star]$, typically}]
\label{lem:syntactic-gap}
If $\Pcal$ contains two distinct programs $p^\star, p^{\star\prime}$ with $\mu_{p^\star}=\mu_{p^{\star\prime}}$, then $\{p^\star\}\subsetneq [p^\star]$.
\end{lemma}

\begin{proof}
Immediate from the definitions.
\end{proof}

\begin{remark}
The hypothesis of Lemma~\ref{lem:syntactic-gap} is automatic in the Solomonoff setting. Any computable environment admits infinitely many syntactic encodings: append a dead code segment, reroute through a universal interpreter, or permute the program's internal representation. No nontrivial semantic class is a singleton in $\Pcal$.
\end{remark}

\begin{theorem}[Strict knowability hierarchy]
\label{thm:strict-hierarchy}
Assume $|\Ocal|\ge 2$, $|\Acal|\ge 2$, and that the hypothesis of Lemma~\ref{lem:syntactic-gap} holds. For any $p^\star\in\Pcal$ and any policy $\pi$ whose support fails at some reachable history $h_{t_0}$ and action $b$ with $\mu_{p^\star}(\cdot\mid h_{t_0},b)$ non-degenerate, and any realized history $h_t$ with $\Prob_{p^\star,\pi}(H_t=h_t)>0$ for which there exists a next action $a$ satisfying $\pi_{t+1}(a\mid h_t)>0$ and $\mu_{p^\star}(\cdot\mid h_t,a)$ non-degenerate,
\[
\{p^\star\}\subsetneq [p^\star]\subsetneq R_\pi(p^\star)\subsetneq R_{h_t}(p^\star)\subseteq \Pcal.
\]
\end{theorem}

\begin{proof}
The three strict inclusions are furnished by Lemma~\ref{lem:syntactic-gap}, Theorem~\ref{thm:policy-gap}, and Lemma~\ref{lem:fit-gap} respectively. Nesting is Lemma~\ref{lem:nesting}.
\end{proof}

\begin{remark}
Table~\ref{tab:hierarchy} records the role of each set in the paper.
\end{remark}

\begin{table}[ht]
\centering
\small
\caption{The four nested sets read as indistinguishability classes at increasing observational power, with the results that witness strict separation.}
\label{tab:hierarchy}
\begin{tabularx}{\textwidth}{@{}>{\raggedright\arraybackslash}p{0.20\textwidth}>{\raggedright\arraybackslash}p{0.34\textwidth}>{\raggedright\arraybackslash}X>{\raggedright\arraybackslash}p{0.16\textwidth}@{}}
\toprule
\textbf{Set} & \textbf{Observer that cannot detect a swap from $p^\star$} & \textbf{Strictly contains} & \textbf{Witnessed by} \\
\midrule
$\{p^\star\}$ & Reads the source code & Nothing & --- \\
$[p^\star]$ & Runs any policy and sees its induced law & $\{p^\star\}$ whenever $[p^\star]$ has $>1$ element & Lemma~\ref{lem:syntactic-gap} \\
$R_\pi(p^\star)$ & Sees the full law of histories under fixed $\pi$ & $[p^\star]$ when $\pi$ leaves some pair unqueried & Theorem~\ref{thm:policy-gap} \\
$R_{h_t}(p^\star)$ & Sees only the realized transcript $h_t$ & $R_\pi(p^\star)$ when some supported next action at $h_t$ is non-degenerate & Lemma~\ref{lem:fit-gap} \\
\bottomrule
\end{tabularx}
\end{table}

The rest of the paper is organized by which set the learner's posterior is trying to concentrate on. The belief side of Section~\ref{sec:belief} delivers concentration on $R_{h_t}(p^\star)$ in a predictive-merging sense. Section~\ref{sec:near-miss} shows that the same Bayes/Gibbs/expected-free-energy architecture need not concentrate beyond that under a substituted finite-support prior: its posterior on $[p^\star]$ may stay bounded away from one. Sections~\ref{sec:semantic-action} and~\ref{sec:main} give the sufficient conditions under which the posterior does concentrate on $[p^\star]$.

\section{The Two-Observer Functional and the Meeting-Point Theorem}
\label{sec:functional}

The problem box of Section~\ref{sec:intro} asks for a triple --- prior, belief rule, action rule --- under which the posterior mass on $[p^\star]$ converges almost surely to one. The four-set hierarchy $\{p^\star\}\subseteq[p^\star]\subseteq R_\pi(p^\star)\subseteq R_{h_t}(p^\star)$ of Section~\ref{sec:hierarchy} identifies the target $[p^\star]$ as the indistinguishability class of the behavioral observer $\omega_{\mathrm{behav}}$ and as the finest set any history-based learner can hope to recover (Theorem~\ref{thm:factor-through-quotient}). This section packages the paper's answer to the box as a single variational object: a functional $\mathcal J_t^{\omega_B,\omega_A}$ indexed by two observers --- one tracking the belief side, one tracking the action side --- whose minimizer is the triple the problem box asks for. The section's punchline is the \emph{meeting-point theorem} (Theorem~\ref{thm:meeting-point}): belief admissibility is a half-lattice above $\omega_{\mathrm{behav}}$ and action admissibility is a half-lattice below, so the two ranges meet at the single observer $\omega_{\mathrm{behav}}$, with the canonical quotient-free belief lift at $\omega_{\mathrm{syn}}$. The two-observer indexing is therefore not a stylistic convenience; it is forced by the refinement lattice of observers that organizes the four-set hierarchy.

The results of this section use only the observer structure of Section~\ref{sec:hierarchy} and the universal prior $w$ of Definition~\ref{def:universal-prior}; the belief-side and action-side constructions they license are identified in Sections~\ref{sec:belief} and~\ref{sec:semantic-action} respectively, and the convergence claim is settled in Section~\ref{sec:main}.

\subsection{Quotient data at an observer}
\label{sec:quotient-data}

Fix an observer $\omega$ with equivalence $\sim_\omega$ on $\Pcal$, and write $\Ccal^\omega:=\Pcal/\!\sim_\omega$ for the set of $\omega$-classes. The prior pushes forward to
\[
\bar w^\omega(c):=\sum_{p\in c}w(p),\qquad c\in\Ccal^\omega,
\]
and the Bayes posterior on $\Pcal$ marginalizes to the class posterior
\[
q_t^\star(c\mid h_t):=\sum_{p\in c}q_t^\star(p\mid h_t).
\]
When $\omega\ge\omega_{\mathrm{behav}}$, every $\omega$-class $c$ is contained in a single behavioral class, so all $p\in c$ share the same interactive law; the class likelihood $\mu_c:=\mu_p$ for any $p\in c$ is then well-defined. When $\omega<\omega_{\mathrm{behav}}$, an $\omega$-class may contain programs with distinct interactive laws, and $\mu_c$ is multi-valued at the class level.

Regardless of where $\omega$ sits relative to $\omega_{\mathrm{behav}}$, the posterior-weighted class-predictive law
\[
\mu_c^{q_t^\star}(o\mid h_t,a):=\frac{1}{q_t^\star(c\mid h_t)}\sum_{p\in c}q_t^\star(p\mid h_t)\,\mu_p(o\mid h_t,a)
\]
is well-defined for every $c$ with $q_t^\star(c\mid h_t)>0$; and so is its complement
\[
\Qcal_t^{-c,\omega}(o\mid h_t,a):=\frac{1}{1-q_t^\star(c\mid h_t)}\sum_{c'\in\Ccal^\omega\setminus\{c\}}q_t^\star(c'\mid h_t)\,\mu_{c'}^{q_t^\star}(o\mid h_t,a).
\]
When $\omega\ge\omega_{\mathrm{behav}}$, $\mu_c^{q_t^\star}=\mu_c$ and the posterior weights in the numerator cancel; we write $\mu_c$ without the posterior superscript in this regime.

\begin{definition}[Bhattacharyya separation at observer $\omega$]
\label{def:bhat-omega}
The \emph{$\omega$-class Bhattacharyya separation} of class $c\in\Ccal^\omega$ under action $a\in\Acal$ at history $h_t$ with $0<q_t^\star(c\mid h_t)<1$ is
\[
B_t^\omega(c,a;h_t):=-\log\sum_{o\in\Ocal}\sqrt{\mu_c^{q_t^\star}(o\mid h_t,a)\,\Qcal_t^{-c,\omega}(o\mid h_t,a)}\;\in\;[0,\infty].
\]
At $\omega=\omega_{\mathrm{behav}}$ we write $B_t:=B_t^{\omega_{\mathrm{behav}}}$; this is the separation used by the class-against-complement action primitive of Section~\ref{sec:semantic-action}.
\end{definition}

Finally, any policy $\pi$ whose action at step $t+1$ is generated by first sampling a class $C\in\Ccal^{\omega_A}$ and then selecting an action targeting that class induces a \emph{class-targeting distribution} $\nu_t^\pi(\cdot\mid h_t)$ on $\Ccal^{\omega_A}$. The class-against-complement action rule of Definition~\ref{def:semantic-rule} fits this template with $\omega_A=\omega_{\mathrm{behav}}$; for now it suffices to know that $\nu_t^\pi$ is well-defined for the policies Section~\ref{sec:semantic-action} will construct.

\subsection{The two-observer functional}
\label{sec:two-observer-functional}

Fix positive scalars $\beta,\gamma>0$ and two observers $\omega_B,\omega_A$.

\begin{definition}[Two-observer variational functional]
\label{def:two-observer-functional}
For a distribution $q$ on $\Pcal$ and a policy $\pi$, define
\[
\boxed{\;
\mathcal J_t^{\omega_B,\omega_A}[q,\pi;h_t]
\;:=\;
\underbrace{\E_q[\ellp_t(p;h_t)]+\KL(q\givena w)}_{\mathcal F_t[q;h_t]}
\;-\;\beta\,\mathcal A_t^{\omega_A}[\pi;h_t,q]
\;+\;\gamma\,\KL\!\bigl(\nu_t^\pi\,\|\,\bar w^{\omega_A}\bigr),
\;}
\]
where $\ellp_t(p;h_t):=-\log\mu_p(o_{1:t}\mid a_{1:t})$ is the history codelength, $\mathcal F_t$ is the algorithmic free energy referenced by the belief observer $\omega_B$ through its pushforward $\bar q^{\omega_B}$ (Section~\ref{sec:belief}), and
\[
\mathcal A_t^{\omega_A}[\pi;h_t,q]:=\E_{C\sim\nu_t^\pi}\bigl[B_t^{\omega_A}(C,A_\pi;h_t)\bigr]
\]
is the expected $\omega_A$-class Bhattacharyya separation under $\pi$'s class-targeting distribution on $\Ccal^{\omega_A}$.
\end{definition}

A parallel \emph{mutual-information variant} replaces the action term:
\[
\mathcal J_t^{\omega_B,\omega_A,\mathrm{MI}}[q,\pi;h_t]
:=\mathcal F_t[q;h_t]-\beta\,\I_q\!\bigl(C^{\omega_A};O\mid A_\pi,h_t\bigr),
\]
where $C^{\omega_A}$ is the random variable coding the $\omega_A$-class of the true program under $q$. At $\omega_A=\omega_{\mathrm{syn}}$, $C^{\omega_A}=P$ and this recovers the standard expected-free-energy information term. The two variants differ by a swap of one R\'enyi divergence for another: the Bhattacharyya action term is $\tfrac{1}{2}D_{1/2}$, the MI term is $D_1=\KL$ under the conventional active-inference decomposition.

\begin{definition}[Meeting-point specialization]
\label{def:meeting-point-shorthand}
Write $\mathcal J_t^{\omega_{\mathrm{behav}}}:=\mathcal J_t^{\omega_{\mathrm{syn}},\omega_{\mathrm{behav}}}$ for the specialization with the canonical belief observer $\omega_B=\omega_{\mathrm{syn}}$ (the quotient-free top of the refinement lattice) and action observer $\omega_A=\omega_{\mathrm{behav}}$ (the meeting point identified in Theorem~\ref{thm:meeting-point}).
\end{definition}

The subsections that follow establish the two-sided admissibility structure of $\mathcal J_t^{\omega_B,\omega_A}$: the belief observer $\omega_B$ must lie above $\omega_{\mathrm{behav}}$ in the refinement lattice, and the action observer $\omega_A$ must lie at or below. These two conditions pick out $\omega_{\mathrm{behav}}$ as the unique single-observer admissible choice and single out the pair $(\omega_{\mathrm{syn}},\omega_{\mathrm{behav}})$ as the canonical two-observer choice.

\subsection{Belief admissibility: a half-lattice above $\omega_{\mathrm{behav}}$}
\label{sec:belief-admissibility}

Let $q_t^{\omega_B}(\cdot\mid h_t)$ denote Bayes applied at the level of $\omega_B$-classes, using the pushforward prior $\bar w^{\omega_B}$ and the class likelihood $\mu_c$ (when well-defined). Two observations bound where $\omega_B$ can sit.

\begin{proposition}[Belief-observer invariance above $\omega_{\mathrm{behav}}$]
\label{prop:belief-invariance-above}
Let $\omega_B$ be any observer with $\omega_B\ge\omega_{\mathrm{behav}}$. Then for every behavioral class $c_{\mathrm{behav}}\in\Pcal/\!\sim_{\mathrm{behav}}$ and every history $h_t$,
\[
\sum_{\substack{c\in\Ccal^{\omega_B}\\ c\subseteq c_{\mathrm{behav}}}}q_t^{\omega_B}(c\mid h_t)\;=\;q_t^{\omega_{\mathrm{behav}}}(c_{\mathrm{behav}}\mid h_t).
\]
In particular, $\omega_B$-level Bayes and $\omega_{\mathrm{behav}}$-level Bayes produce identical $\omega_{\mathrm{behav}}$-class posterior dynamics.
\end{proposition}

\begin{proof}
Because $\omega_B\ge\omega_{\mathrm{behav}}$, every $\omega_B$-class $c$ lies inside a single behavioral class, so $\mu_c=\mu_{c_{\mathrm{behav}}}$ for each $c\subseteq c_{\mathrm{behav}}$. Direct computation:
\begin{align*}
\sum_{c\subseteq c_{\mathrm{behav}}}q_t^{\omega_B}(c\mid h_t)
&=\sum_{c\subseteq c_{\mathrm{behav}}}\frac{\bar w^{\omega_B}(c)\,\mu_c(o_{1:t}\mid a_{1:t})}{\barxi(o_{1:t}\mid a_{1:t})}\\
&=\frac{\mu_{c_{\mathrm{behav}}}(o_{1:t}\mid a_{1:t})}{\barxi(o_{1:t}\mid a_{1:t})}\sum_{c\subseteq c_{\mathrm{behav}}}\bar w^{\omega_B}(c)
=\frac{\bar w^{\omega_{\mathrm{behav}}}(c_{\mathrm{behav}})\,\mu_{c_{\mathrm{behav}}}(o_{1:t}\mid a_{1:t})}{\barxi(o_{1:t}\mid a_{1:t})},
\end{align*}
which equals $q_t^{\omega_{\mathrm{behav}}}(c_{\mathrm{behav}}\mid h_t)$. The $\omega_B$-level and program-level mixtures coincide with $\barxi$ because pushforward preserves the mixture identity.
\end{proof}

\begin{proposition}[Belief-observer ill-typing below $\omega_{\mathrm{behav}}$]
\label{prop:belief-illtyped-below}
If $\omega_B<\omega_{\mathrm{behav}}$, there exists an $\omega_B$-class $c\in\Ccal^{\omega_B}$ containing two programs $p_1,p_2\in c$ with $\mu_{p_1}\neq\mu_{p_2}$. Consequently the class likelihood $\mu_c$ is not well-defined as a single distribution, and $\omega_B$-level Bayes is ill-typed without auxiliary within-class sub-posterior structure.
\end{proposition}

\begin{proof}
$\omega_B<\omega_{\mathrm{behav}}$ means $\omega_{\mathrm{behav}}\not\le\omega_B$ in the refinement preorder (Section~\ref{sec:hierarchy}), i.e., $p_1\sim_{\omega_B}p_2$ does not imply $p_1\sim_{\mathrm{behav}}p_2$. Pick witnesses $p_1,p_2$ with $p_1\sim_{\omega_B}p_2$ but $\mu_{p_1}\neq\mu_{p_2}$; they belong to the same $\omega_B$-class but induce distinct interactive laws. The assignment $c\mapsto\mu_c$ is then multi-valued on $c=[p_1]_{\omega_B}$.
\end{proof}

Proposition~\ref{prop:belief-invariance-above} says: above the behavioral observer, \emph{any} belief observer produces identical $\omega_{\mathrm{behav}}$-class dynamics, so the choice $\omega_B\in[\omega_{\mathrm{behav}},\omega_{\mathrm{syn}}]$ is a lift rather than a real parameter. Proposition~\ref{prop:belief-illtyped-below} says: below the behavioral observer, class-level Bayes is not a single procedure. Belief admissibility is therefore the half-lattice $\{\omega:\omega\ge\omega_{\mathrm{behav}}\}$, closed at $\omega_{\mathrm{behav}}$ and open-ended above.

\subsection{Action admissibility: a half-lattice below $\omega_{\mathrm{behav}}$}
\label{sec:action-admissibility}

The symmetric bound comes from the action side. Here the constraint is that separating two $\omega_A$-classes demands classes with distinct laws; above $\omega_{\mathrm{behav}}$, the lattice delivers pairs of classes that are behaviorally indistinguishable.

\begin{proposition}[Action-observer cap]
\label{prop:action-cap}
Let $\omega_A$ be an observer with $\omega_A>\omega_{\mathrm{behav}}$, and let $c_1,c_2\in\Ccal^{\omega_A}$ be two distinct $\omega_A$-classes contained in a single behavioral class, with $q_t^\star(c_1\mid h_t),q_t^\star(c_2\mid h_t)>0$. Set $\alpha:=q_t^\star(c_2\mid h_t)/(1-q_t^\star(c_1\mid h_t))\in(0,1]$. Then for every action $a\in\Acal$,
\[
B_t^{\omega_A}(c_1,a;h_t)\;\le\;\tfrac{1}{2}\log(1/\alpha)
\;=\;\tfrac{1}{2}\log\!\frac{1-q_t^\star(c_1\mid h_t)}{q_t^\star(c_2\mid h_t)},
\]
a bound independent of $a$.
\end{proposition}

\begin{proof}
Since $c_1,c_2$ share a behavioral class, $\mu_{c_1}(o\mid h_t,a)=\mu_{c_2}(o\mid h_t,a)=:\mu(o)$ for every $o$ and $a$. Writing $\Rcal$ for the posterior mixture over $\omega_A$-classes other than $c_1$ and $c_2$,
\[
\Qcal_t^{-c_1,\omega_A}(o\mid h_t,a)=\alpha\,\mu(o)+(1-\alpha)\,\Rcal(o).
\]
The pointwise bound $\sqrt{\alpha\mu+(1-\alpha)\Rcal}\ge\sqrt{\alpha\mu}$ gives
\[
\sum_o\sqrt{\mu_{c_1}(o)\,\Qcal_t^{-c_1,\omega_A}(o)}=\sum_o\sqrt{\mu(o)\bigl(\alpha\mu(o)+(1-\alpha)\Rcal(o)\bigr)}\ge\sqrt{\alpha}\sum_o\mu(o)=\sqrt{\alpha}.
\]
Taking $-\log$ yields $B_t^{\omega_A}(c_1,a;h_t)\le\tfrac{1}{2}\log(1/\alpha)$, independent of $a$.
\end{proof}

\begin{corollary}[Behavioral twins have frozen posterior ratio]
\label{cor:twins-frozen-ratio}
Let $\omega_A>\omega_{\mathrm{behav}}$ and let $c_1,c_2\in\Ccal^{\omega_A}$ be distinct $\omega_A$-classes contained in a single behavioral class. Then for every $t$ and every history $h_t$ with $q_t^\star(c_1\mid h_t),q_t^\star(c_2\mid h_t)>0$,
\[
\frac{q_t^\star(c_1\mid h_t)}{q_t^\star(c_2\mid h_t)}=\frac{\bar w^{\omega_A}(c_1)}{\bar w^{\omega_A}(c_2)}.
\]
In particular, no $\omega_A$-class posterior $q_t^\star(c_i\mid H_t)$ can tend to one almost surely whenever a behavioral twin has positive prior mass.
\end{corollary}

\begin{proof}
Bayes for classes with identical likelihoods $\mu_{c_1}=\mu_{c_2}$ gives
\[
q_t^\star(c_i\mid h_t)=\frac{\bar w^{\omega_A}(c_i)\,\mu_{c_i}(o_{1:t}\mid a_{1:t})}{\barxi(o_{1:t}\mid a_{1:t})},
\]
and the ratio at $c_1,c_2$ reduces to $\bar w^{\omega_A}(c_1)/\bar w^{\omega_A}(c_2)$ at every history.
\end{proof}

The cap theorem says: above the behavioral observer, no action can separate behavioral twins, so the Bhattacharyya term of $\mathcal J_t^{\omega_B,\omega_A}$ cannot drive their posteriors apart; Corollary~\ref{cor:twins-frozen-ratio} is the sharpened version --- the ratio is pinned to the prior ratio, uniformly in $t$. Action admissibility is therefore the symmetric half-lattice $\{\omega:\omega\le\omega_{\mathrm{behav}}\}$, closed at $\omega_{\mathrm{behav}}$ and open-ended below.

\subsection{The meeting-point theorem}
\label{sec:meeting-point}

\begin{theorem}[Meeting point]
\label{thm:meeting-point}
In the $\mathcal J_t^{\omega_B,\omega_A}$ framework, belief admissibility and action admissibility are half-lattices bounded on opposite sides of the refinement order, and they meet at a single observer:
\begin{itemize}
\item[(a)] \textbf{Belief lower bound.} $\omega_B\ge\omega_{\mathrm{behav}}$ is required for $\omega_B$-level Bayes to be well-defined as a single procedure on $\omega_B$-classes (Proposition~\ref{prop:belief-illtyped-below}).
\item[(b)] \textbf{Belief upper freedom.} For every $\omega_B\ge\omega_{\mathrm{behav}}$, $\omega_B$-level Bayes produces identical $\omega_{\mathrm{behav}}$-class posterior trajectories to $\omega_{\mathrm{behav}}$-level Bayes (Proposition~\ref{prop:belief-invariance-above}).
\item[(c)] \textbf{Action upper bound.} $\omega_A\le\omega_{\mathrm{behav}}$ is required for the Bhattacharyya action term $\mathcal A_t^{\omega_A}$ to separate distinct $\omega_A$-classes uniformly in the policy (Proposition~\ref{prop:action-cap} and Corollary~\ref{cor:twins-frozen-ratio}).
\item[(d)] \textbf{Action lower freedom.} For every $\omega_A\le\omega_{\mathrm{behav}}$, $\omega_A$-level action minimization of $\mathcal J_t^{\omega_B,\omega_A}$ supports almost-sure convergence of $q_t^\star([p^\star]_{\omega_A}\mid H_t)\to 1$ under the appropriate separation and rule conditions (Proposition~\ref{prop:goal-dialed}).
\end{itemize}
Consequently $\omega_{\mathrm{behav}}$ is the unique observer at which belief admissibility meets action admissibility: the minimum admissible belief observer equals the maximum useful action observer.
\end{theorem}

\begin{proof}
(a) and (b) are Propositions~\ref{prop:belief-illtyped-below} and~\ref{prop:belief-invariance-above}; (c) is Proposition~\ref{prop:action-cap} specialized to any $\omega_A>\omega_{\mathrm{behav}}$, which by Corollary~\ref{cor:twins-frozen-ratio} freezes the posterior ratio of any pair of behavioral twins independently of the policy; (d) is Proposition~\ref{prop:goal-dialed} below. The belief-admissibility range $\{\omega:\omega\ge\omega_{\mathrm{behav}}\}$ and the action-admissibility range $\{\omega:\omega\le\omega_{\mathrm{behav}}\}$ are the two principal half-lattices of the refinement order bounded at $\omega_{\mathrm{behav}}$; their intersection is the singleton $\{\omega_{\mathrm{behav}}\}$ and uniqueness is the lattice identity $\{\omega_{\mathrm{behav}}\}=\{\omega\ge\omega_{\mathrm{behav}}\}\cap\{\omega\le\omega_{\mathrm{behav}}\}$.
\end{proof}

\begin{corollary}[Canonical two-observer choice]
\label{cor:canonical-pair}
If the framework requires a single-observer specification, $\omega_{\mathrm{behav}}$ is forced. If the framework allows distinct belief and action observers, the canonical pair is $(\omega_{\mathrm{syn}},\omega_{\mathrm{behav}})$: $\omega_{\mathrm{syn}}$ is the quotient-free top of the refinement lattice and $\omega_{\mathrm{behav}}$ is the meeting point. By Proposition~\ref{prop:belief-invariance-above}, the family $\{\mathcal J_t^{\omega_B,\omega_{\mathrm{behav}}}:\omega_B\ge\omega_{\mathrm{behav}}\}$ produces identical $\omega_{\mathrm{behav}}$-class dynamics; $\mathcal J_t^{\omega_{\mathrm{behav}}}=\mathcal J_t^{\omega_{\mathrm{syn}},\omega_{\mathrm{behav}}}$ of Definition~\ref{def:meeting-point-shorthand} is its canonical representative.
\end{corollary}

The meeting-point theorem is the structural content of the two-observer design. The functional's two observer indices are not stylistic: the refinement lattice has a single fixed point where the belief-side admissibility range (closed above $\omega_{\mathrm{behav}}$) meets the action-side usefulness range (closed below $\omega_{\mathrm{behav}}$), and the observer indices record where each half of the learning problem must live. The object $\omega_{\mathrm{behav}}$ is a lattice-theoretic invariant of the paper's framework, not a design choice.

\begin{remark}[Leibniz identification]
\label{rem:leibniz}
Two programs $p_1,p_2$ satisfy $p_1\sim_{\mathrm{behav}}p_2$ iff they are indiscernible under every interactive probing --- the interactive-learning analogue of Leibniz's identity of indiscernibles on $\Pcal$. In these terms, $[p^\star]=[p^\star]_{\mathrm{behav}}$ is the Leibniz invariant of $p^\star$, and Theorem~\ref{thm:meeting-point} identifies this invariant as a lattice fixed point rather than a pragmatic stipulation: the two admissibility ranges already force it before any external criterion of ``behavioral equivalence'' is imposed.
\end{remark}

\subsection{Goal-dialed convergence}
\label{sec:goal-dialed}

The meeting-point theorem fixes a canonical observer; the family $\{\mathcal J_t^{\omega_B,\omega_A}:\omega_A\le\omega_{\mathrm{behav}}\}$ nevertheless has content as a \emph{dial}. Each $\omega_A\le\omega_{\mathrm{behav}}$ corresponds to a coarser learning target $[p^\star]_{\omega_A}\supseteq[p^\star]$, and the convergence argument of Section~\ref{sec:main} specializes to the dialed-down target with minimal change.

\begin{proposition}[Goal-dialed convergence]
\label{prop:goal-dialed}
Fix an observer $\omega_A$ with $\omega_A\le\omega_{\mathrm{behav}}$, and let $[p^\star]_{\omega_A}:=\{p\in\Pcal:p\sim_{\omega_A}p^\star\}$. Suppose:
\begin{itemize}
\item[\textup{(C1$^{\omega_A}$)}] the universal prior of Definition~\ref{def:universal-prior} is in force, so $\bar w^{\omega_A}(c)>0$ for every $c\in\Ccal^{\omega_A}$;
\item[\textup{(C2)}] belief is the Bayes/Gibbs posterior (Lemma~\ref{lem:variational});
\item[\textup{(C3$^{\omega_A}$)}] for every $c\in\Ccal^{\omega_A}$ with $0<q_t^\star(c\mid h_t)<1$, $\sup_{a\in\Acal}B_t^{\omega_A}(c,a;h_t)\ge\eta^{\omega_A}(c)>0$;
\item[\textup{(C4$^{\omega_A}$)}] action follows the class-against-complement rule of Section~\ref{sec:semantic-action} with $\Ccal$ replaced by $\Ccal^{\omega_A}$ throughout.
\end{itemize}
Then $q_t^\star([p^\star]_{\omega_A}\mid H_t)\to 1$ almost surely.
\end{proposition}

\begin{proof}[Proof sketch]
The main-theorem argument of Section~\ref{sec:main} (Theorem~\ref{thm:semantic-convergence}) uses five ingredients: positivity of the pushforward prior on each class; the separation condition (C3); a conditional Borel--Cantelli argument; a Bhattacharyya contraction lemma; and a promotion-supporting lower bound on class targeting. Ingredient~1 transfers by (C1$^{\omega_A}$), since the universal prior assigns positive mass to every computable program and each $\omega_A$-class contains at least one. Ingredient~2 is (C3$^{\omega_A}$). Ingredients~3 and~4 are observer-agnostic: Borel--Cantelli does not depend on the $\omega$ index, and the Bhattacharyya contraction uses only the R\'enyi-$1/2$ identity in the likelihood-ratio calculus. Ingredient~5 follows from Proposition~\ref{prop:semantic-is-promotion-supporting} applied with $\Ccal^{\omega_A}$ in place of $\Ccal$. The place where the choice $\omega_A=\omega_{\mathrm{behav}}$ was load-bearing in Theorem~\ref{thm:semantic-convergence} was the implicit exclusion of $\omega_A>\omega_{\mathrm{behav}}$ (Proposition~\ref{prop:action-cap}); under $\omega_A\le\omega_{\mathrm{behav}}$ no cap applies. Details of the transfer are given alongside Theorem~\ref{thm:semantic-convergence}.
\end{proof}

Proposition~\ref{prop:goal-dialed} reads the family $\{\mathcal J_t^{\omega_B,\omega_A}:\omega_A\le\omega_{\mathrm{behav}}\}$ as an observer dial: minimizing $\mathcal J_t^{\omega_B,\omega_A}$ drives the posterior onto the $\omega_A$-class of $p^\star$, with $\omega_A=\omega_{\mathrm{behav}}$ the finest allowable setting (the knowability ceiling of Theorem~\ref{thm:factor-through-quotient}) and coarser $\omega_A$'s targeting coarser properties of programs.

\subsection{How the rest of the paper reads under the functional}
\label{sec:functional-forward-plan}

The remainder of the paper identifies the triple the problem box asks for as the minimizer of $\mathcal J_t^{\omega_{\mathrm{behav}}}$, then validates the convergence conclusion.

Section~\ref{sec:belief} identifies the \emph{belief-side minimizer} of $\mathcal J_t^{\omega_B,\omega_A}$ at the canonical belief observer $\omega_B=\omega_{\mathrm{syn}}$: the Bayes/Gibbs posterior under the universal prior. The belief-side story is the content of Lemma~\ref{lem:variational} and Lemma~\ref{lem:kl-necessity}, re-read as the $q$-minimization of $\mathcal F_t$.

Section~\ref{sec:semantic-action} identifies the \emph{action-side minimizer} at the canonical action observer $\omega_A=\omega_{\mathrm{behav}}$: the class-against-complement action rule. The class-targeting distribution $\nu_t^\pi$ is fixed explicitly by Definition~\ref{def:semantic-rule}, and Proposition~\ref{prop:semantic-is-promotion-supporting} supplies the promotion-supporting lower bound used in the convergence argument.

Section~\ref{sec:main} establishes the \emph{convergence payoff}: under the universal prior, the Bayes/Gibbs belief, the class-against-complement rule, and the semantic separation condition (C3), the posterior mass on $[p^\star]$ converges almost surely to one, which is $\mathcal J_t^{\omega_{\mathrm{behav}}}$-minimization delivering the answer the problem box asks for.

Section~\ref{sec:near-miss} reads the algorithmic-free-energy principle as the \emph{wrong-observer} instance of the same functional: its architecture coincides with the minimizer of the MI variant $\mathcal J_t^{\omega_{\mathrm{syn}},\mathrm{MI}}$ at action observer $\omega_A=\omega_{\mathrm{syn}}$, which by Proposition~\ref{prop:action-cap} lies above the meeting point and therefore cannot separate behavioral twins; the near-miss is the boundary example outside the admissible range, not a defeat of Bayesian updating.

\section{Belief at $\omega_{\mathrm{syn}}$: Universal Bayesian Belief}
\label{sec:belief}

Section~\ref{sec:functional} fixed the canonical belief observer at the top of the refinement lattice: $\omega_B=\omega_{\mathrm{syn}}$ (Corollary~\ref{cor:canonical-pair}). By Proposition~\ref{prop:belief-invariance-above}, this choice is invariant --- any $\omega_B\in[\omega_{\mathrm{behav}},\omega_{\mathrm{syn}}]$ produces the same $\omega_{\mathrm{behav}}$-class posterior dynamics --- so nothing is lost by working at the quotient-free top. What is gained is a clean setting in which to identify the belief-side minimizer of $\mathcal J_t^{\omega_{\mathrm{syn}},\omega_A}$: Bayes/Gibbs on programs, under the universal interactive prior, with its algorithmic free energy as the belief term.

This section makes that object explicit. Definition~\ref{def:universal-prior} fixes the universal prior of Solomonoff and Hutter. Definition~\ref{def:afe} identifies the \emph{algorithmic free energy} $\Fcal_t$ with the belief term of $\mathcal J_t^{\omega_{\mathrm{syn}},\omega_A}$ at the canonical belief observer; Lemma~\ref{lem:variational} shows that its unique minimizer over $q$ is the exact Bayes/Gibbs posterior; Lemma~\ref{lem:kl-necessity} shows the KL-based form is forced within the natural convex class; Lemma~\ref{lem:merging} shows that the resulting posterior predictive law merges almost surely into the history-compatible set $R_{h_t}(p^\star)$ --- the outermost of the four sets of Section~\ref{sec:hierarchy}. The belief side alone therefore fixes the posterior at the coarsest observer in the four-set hierarchy, $\omega_{h_t}$; the gap from $R_{h_t}(p^\star)$ down to $[p^\star]$ is what the action side of Sections~\ref{sec:semantic-action} and~\ref{sec:main} closes.

\subsection{Universal prior}

\begin{definition}[Universal interactive prior]
\label{def:universal-prior}
For $p\in\Pcal$, let $\ellp(p)=|p|$ be the self-delimiting code length. Define raw weights $\wt w(p)=2^{-\ellp(p)}$, normalizing constant $Z=\sum_{p\in\Pcal}\wt w(p)\le 1$, normalized prior
\[
w(p)=\frac{\wt w(p)}{Z},
\]
and universal interactive mixture
\[
\barxi(o_{1:t}\givena a_{1:t})
:=
\sum_{p\in\Pcal} w(p)\mu_p(o_{1:t}\givena a_{1:t}).
\]
\end{definition}

\begin{lemma}[Machine invariance]
\label{lem:prior-invariance}
Let $U$ and $V$ be universal prefix machines, with corresponding normalized universal interactive mixtures $\barxi_U$ and $\barxi_V$. Then there are constants $c_{U,V}, c_{V,U}>0$ such that
\[
\barxi_U(o_{1:t}\givena a_{1:t})\ge c_{U,V}\barxi_V(o_{1:t}\givena a_{1:t}),
\qquad
\barxi_V(o_{1:t}\givena a_{1:t})\ge c_{V,U}\barxi_U(o_{1:t}\givena a_{1:t})
\]
for every $t$ and every compatible interaction prefix. Consequently the semantic complexity $K_U(\nu):=\min\{|p|:\mu_p=\nu\}$ of any computable interactive law $\nu$ is invariant up to an additive constant across universal machines.
\end{lemma}

\begin{proof}
Let $Z_U:=\sum_{q\in\Pcal_U}2^{-|q|}$ and $Z_V:=\sum_{p\in\Pcal_V}2^{-|p|}$. Because $U$ is universal, there is a fixed self-delimiting translator $\tau_{V\to U}$ such that for every $p\in\Pcal_V$, the concatenated $U$-program $\tau_{V\to U}p$ lies in $\Pcal_U$ and induces the same interactive law. Hence
\[
\barxi_U(o_{1:t}\givena a_{1:t})
\ge
\frac{1}{Z_U}\sum_{p\in\Pcal_V}2^{-|\tau_{V\to U}p|}\mu_p^V(o_{1:t}\givena a_{1:t})
=
2^{-|\tau_{V\to U}|}\frac{Z_V}{Z_U}\barxi_V(o_{1:t}\givena a_{1:t}).
\]
The reverse inequality is symmetric. For semantic complexity, if $p_V$ is a shortest $V$-program realizing $\nu$, then $\tau_{V\to U}p_V$ is a $U$-program realizing $\nu$, so $K_U(\nu)\le K_V(\nu)+|\tau_{V\to U}|$. Swapping $U$ and $V$ gives $K_V(\nu)\le K_U(\nu)+|\tau_{U\to V}|$, hence additive invariance.
\end{proof}

\begin{lemma}[{Positive prior mass on $[p^\star]$ is necessary}]
\label{lem:prior-necessity}
Let $w'$ be any prior on $\Pcal$, with induced mixture $\xi'(o_{1:t}\givena a_{1:t}):=\sum_p w'(p)\mu_p(o_{1:t}\givena a_{1:t})$. If $\xi'$ fails to dominate some computable interactive law $\nu$, and $c_\nu\in\Ccal$ is its semantic class, then $w'(c_\nu)=0$, and along any history generated by $\nu$ the posterior under $w'$ satisfies $q_t^{w'}(c_\nu\mid h_t)=0$ whenever defined.
\end{lemma}

\begin{proof}
If $w'(c_\nu)>0$, there exists $p\in c_\nu$ with $w'(p)>0$, and $\xi'(o_{1:t}\givena a_{1:t})\ge w'(p)\nu(o_{1:t}\givena a_{1:t})$, so $\xi'$ dominates $\nu$. Contradiction. Posterior mass on $c_\nu$ is a numerator sum over a measure-zero set, so it is zero.
\end{proof}

\subsection{Algorithmic free energy as the belief term of $\mathcal J_t$ at $\omega_B=\omega_{\mathrm{syn}}$}

For $p\in\Pcal$ and history $h_t$, define the cumulative history-fit loss
\[
L_t(p;h_t):=-\log\mu_p(o_{1:t}\givena a_{1:t}),
\]
with the convention $L_t=+\infty$ if the likelihood vanishes. Consistent with the notation of Definition~\ref{def:two-observer-functional}, $L_t(p;h_t)=\ellp_t(p;h_t)$.

\begin{definition}[Algorithmic free energy]
\label{def:afe}
For a posterior $q\in\Delta(\Pcal)$ with $\E_q[L_t]<\infty$ and $\KL(q\|w)<\infty$, the \emph{algorithmic free energy} at history $h_t$ is
\[
\Fcal_t[q;h_t]:=\E_q[L_t(p;h_t)]+\KL(q\|w).
\]
By construction, $\Fcal_t[q;h_t]$ is the belief term of the two-observer functional $\mathcal J_t^{\omega_{\mathrm{syn}},\omega_A}[q,\pi;h_t]$ at the canonical belief observer $\omega_B=\omega_{\mathrm{syn}}$ (Definition~\ref{def:two-observer-functional}): it is the part that depends on $q$ alone and not on $\pi$.
\end{definition}

\begin{lemma}[Variational identity]
\label{lem:variational}
Assume $\barxi(o_{1:t}\givena a_{1:t})>0$, and define the Bayes/Gibbs posterior
\[
q_t^\star(p\mid h_t):=\frac{w(p)\mu_p(o_{1:t}\givena a_{1:t})}{\barxi(o_{1:t}\givena a_{1:t})}.
\]
Then for every admissible posterior $q$,
\[
\Fcal_t[q;h_t]
=
\KL\!\bigl(q\,\|\,q_t^\star(\cdot\mid h_t)\bigr)-\log\barxi(o_{1:t}\givena a_{1:t}),
\]
and $\Fcal_t$ is uniquely minimized at $q_t^\star$.
\end{lemma}

\begin{proof}
By definition,
\[
\Fcal_t[q;h_t]=\sum_p q(p)\log\frac{q(p)}{w(p)\mu_p(o_{1:t}\givena a_{1:t})}
=\KL\!\bigl(q\,\|\,q_t^\star(\cdot\mid h_t)\bigr)-\log\barxi(o_{1:t}\givena a_{1:t}).
\]
Uniqueness of the minimizer follows from strict positivity of the KL divergence away from $q_t^\star$.
\end{proof}

\begin{remark}[Belief-side minimizer of $\mathcal J_t^{\omega_{\mathrm{syn}},\omega_A}$]
\label{rem:afe-bayes-unification}
Because $\Fcal_t$ is the $q$-dependent part of $\mathcal J_t^{\omega_{\mathrm{syn}},\omega_A}[q,\pi;h_t]$, Lemma~\ref{lem:variational} identifies the belief-side minimizer of the two-observer functional at $\omega_B=\omega_{\mathrm{syn}}$ with the exact Bayes/Gibbs posterior $q_t^\star$, uniformly in the action-side choice of $\pi$ and $\omega_A$. Put differently, fixing $\omega_B=\omega_{\mathrm{syn}}$ turns the $q$-minimization of $\mathcal J_t^{\omega_B,\omega_A}$ into Bayes under the universal interactive prior; no approximation is involved and no action-side choice disturbs this identification.

The algorithmic free energy of Definition~\ref{def:afe} has not, to our knowledge, appeared in this exact form. The variational free energy of active inference takes a Gibbs form but uses a stipulated parametric prior that upper-bounds model evidence only approximately \citep{Friston2010,ParrFriston2019}; Solomonoff--Hutter Bayesian prediction uses the Kolmogorov-complexity-weighted universal prior but does not state a free-energy functional \citep{Solomonoff1964a,Solomonoff1964b,Hutter2005}. Coupling the two --- the Gibbs form with the universal prior as reference measure --- collapses the distinction. Under the universal prior the algorithmic free energy is not a bound on the evidence, it is the negative log-evidence up to an additive constant, and its minimizer is the exact Bayes/Gibbs posterior. The variational structure of active inference is preserved; the stipulative quality of the prior is removed.
\end{remark}

\begin{lemma}[KL is forced by the exact Gibbs variational formula]
\label{lem:kl-necessity}
Let $D(\cdot\|w)$ be a proper, convex, lower-semicontinuous functional on $\ell^1(\Pcal)$, extended by $+\infty$ off $\Delta(\Pcal)$. Assume that for every bounded loss profile $L\in\ell^\infty(\Pcal)$,
\[
\inf_{q\in\Delta(\Pcal)}\bigl\{\E_q[L]+D(q\|w)\bigr\}
=
-\log\sum_{p\in\Pcal} w(p)e^{-L(p)}.
\]
Then $D(q\|w)=\KL(q\|w)$ for every $q\in\Delta(\Pcal)$.
\end{lemma}

\begin{proof}
For bounded $f\in\ell^\infty(\Pcal)$ define $H(f):=\log\sum_p w(p)e^{f(p)}$. With $f=-L$, the hypothesis is equivalent to $H(f)=\sup_q\{\E_q[f]-D(q\|w)\}$. Fenchel--Moreau on the pair $(\ell^1,\ell^\infty)$ gives $D(q\|w)=\sup_f\{\E_q[f]-H(f)\}$.

If $\KL(q\|w)=\infty$, the upper bound $D(q\|w)\le \KL(q\|w)$ is immediate. Otherwise $q\ll w$. For $q_f(p):=w(p)e^{f(p)}/\sum_r w(r)e^{f(r)}$,
\[
\E_q[f]-H(f)=\KL(q\|w)-\KL(q\|q_f)\le \KL(q\|w),
\]
so $D(q\|w)\le\KL(q\|w)$. For the reverse, first suppose $q\not\ll w$, and let $A:=\{p\in\Pcal:q(p)>0,\ w(p)=0\}$. Then $q(A)>0$. For $f_n:=n\,\one_A$, one has $H(f_n)=\log\sum_p w(p)e^{f_n(p)}=0$ and $\E_q[f_n]=nq(A)\to\infty$, so $D(q\|w)=\infty=\KL(q\|w)$. It remains to treat the case $q\ll w$. Set $r(p):=q(p)/w(p)$ on $\{w(p)>0\}$ and $f_n(p):=\max\{-n,\min\{\log r(p),n\}\}$. Each $f_n$ is bounded. Dominated convergence gives $H(f_n)\to 0$. If $\KL(q\|w)<\infty$, $f_n(p)\to\log r(p)$ pointwise; the positive part is $q$-integrable by finiteness of $\KL$, and the bound $\sum_{r(p)<1}w(p)(-r\log r)\le 1/e$ controls the negative part, so $\E_q[f_n]\to\KL(q\|w)$. If $\KL=\infty$, monotone convergence on $f_n^+$ gives $\E_q[f_n^+]\to\infty$, and the same bound on the negative part shows $\sup_n\E_q[f_n^-]<\infty$, hence $\E_q[f_n]\to\infty$. Either way $D(q\|w)\ge\sup_n\{\E_q[f_n]-H(f_n)\}=\KL(q\|w)$.
\end{proof}

\begin{remark}
Lemma~\ref{lem:kl-necessity} singles out the KL-based free energy within the natural convex class.
\end{remark}

\subsection{Predictive merging into $R_{h_t}(p^\star)$}

\begin{lemma}[Predictive merging]
\label{lem:merging}
If the interaction is generated by $p^\star$ under any history-adapted policy, then under the universal prior and Bayes/Gibbs posterior
\[
\sum_{t=0}^{\infty}
\KL\!\bigl(\mu_{p^\star}(\cdot\mid h_t,a_{t+1})\,\|\,\Qcal_t^O(\cdot\mid h_t,a_{t+1})\bigr)
<\infty
\qquad
\text{almost surely,}
\]
where
\[
\Qcal_t^O(o\mid h_t,a):=\sum_{p\in\Pcal} q_t^\star(p\mid h_t)\mu_p(o\mid h_t,a).
\]
In particular the one-step predictive KL tends to zero almost surely.
\end{lemma}

\begin{proof}
Mixture dominance gives $\barxi(o_{1:T}\givena a_{1:T})\ge w(p^\star)\mu_{p^\star}(o_{1:T}\givena a_{1:T})$, so $\log(\mu_{p^\star}/\barxi)\le -\log w(p^\star)$. Sequential factoring and conditional expectation under the true law yield
\[
\sum_{t=0}^{T-1}\E\!\left[\KL\!\bigl(\mu_{p^\star}(\cdot\mid h_t,a_{t+1})\,\|\,\Qcal_t^O(\cdot\mid h_t,a_{t+1})\bigr)\right]
\le -\log w(p^\star).
\]
Monotone convergence gives finiteness of the infinite sum almost surely, hence summand convergence to zero.
\end{proof}

\begin{remark}[Belief alone reaches the outermost set]
\label{rem:belief-reaches-Rht}
Predictive merging places the Bayes/Gibbs posterior, in a one-step predictive sense, inside the outermost set $R_{h_t}(p^\star)$ of the four-set hierarchy of Section~\ref{sec:hierarchy}: the posterior predictive law matches the true next-step law along the realized interaction. This is the Solomonoff-interactive analogue of the classical merging-of-opinions phenomenon \citep{BlackwellDubins1962}. In observer terms, the belief-side minimizer of $\mathcal J_t^{\omega_{\mathrm{syn}},\omega_A}$ under the universal prior reaches the history observer $\omega_{h_t}$; the two remaining strict inclusions $R_{h_t}(p^\star)\supsetneq R_\pi(p^\star)\supsetneq[p^\star]$ of Theorem~\ref{thm:strict-hierarchy} are not closed by belief alone. Closing them requires an action-side minimizer that lifts the realized observer past $\omega_\pi$ toward $\omega_{\mathrm{behav}}$ --- the content of Sections~\ref{sec:semantic-action} and~\ref{sec:main}.
\end{remark}

\section{The Algorithmic Free Energy Principle and Its Near-Miss}
\label{sec:near-miss}

The belief side of Section~\ref{sec:belief} settles one ingredient. The remaining question is action: does the learner's realized policy induce an observer that refines toward the behavioral observer $\omega_{\mathrm{behav}}$, thereby closing the strict gap $\omega_\pi<\omega_{\mathrm{behav}}$ of Theorem~\ref{thm:policy-gap}? The natural generic answer, expected free energy, is a near-miss.

\subsection{Expected free energy}

\begin{definition}[Expected free energy]
\label{def:efe}
At history $h_t$, define the joint predictive law over latent program $P$ and next observation $O$ by
\[
\Qcal_t^{PO}(p,o\mid h_t,a):=q_t^\star(p\mid h_t)\mu_p(o\mid h_t,a).
\]
With observation marginal $\Qcal_t^O(\cdot\mid h_t,a)$, preference distribution $\rho_t(\cdot\mid h_t,a)$, and action cost $c_t(a;h_t)$, the \emph{expected free energy} is
\[
\Gcal_t(a;h_t)
:=
c_t(a;h_t)
+
\E_{\Qcal_t^{PO}}\!\bigl[\log q_t^\star(P\mid h_t)-\log q_{t+1}^\star(P\mid h_t,a,O)-\log\rho_t(O\mid h_t,a)\bigr].
\]
\end{definition}

\begin{lemma}[Risk minus information gain]
\label{lem:risk-ig}
For every history $h_t$ and action $a$,
\[
\Gcal_t(a;h_t)
=
c_t(a;h_t)
+\E_{o\sim\Qcal_t^O}\bigl[-\log\rho_t(o\mid h_t,a)\bigr]
-\I_{\Qcal_t^{PO}}(P;O\mid h_t,a).
\]
\end{lemma}

\begin{proof}
The first expectation depends only on the observation marginal and equals the risk term. The second expectation equals $-\I_{\Qcal_t^{PO}}(P;O\mid h_t,a)$ by the standard mutual-information identity. Summing gives the claim.
\end{proof}

\begin{corollary}[Expected free energy as a unifying epistemic principle]
\label{cor:efe-specialization}
Fix zero action cost and uniform preferences $\rho_t$ on a common finite feasible observation set. Then expected free energy minimization is equivalent, up to an action-independent constant, to maximization of the mutual information $\I_{\Qcal_t^{PO}}(P;O\mid h_t,a)$ between the latent program and the next observation.
\end{corollary}

\begin{proof}
By Lemma~\ref{lem:risk-ig}, $\Gcal_t(a;h_t)=c_t(a;h_t)+\E_{o\sim\Qcal_t^O}[-\log\rho_t(o\mid h_t,a)]-\I_{\Qcal_t^{PO}}(P;O\mid h_t,a)$. Zero action cost annihilates the first term. Uniform $\rho_t$ on a common finite feasible set makes the second term action-independent. The displayed specialization follows.
\end{proof}

\begin{remark}
\label{rem:efe-unification}
Corollary~\ref{cor:efe-specialization} is the action-side unification. In this regime, expected-free-energy minimization is simultaneously the Bayesian experimental design criterion \citep{Lindley1956,ChalonerVerdinelli1995,MacKay1992}, the compression-based curiosity criterion \citep{Schmidhuber1991,Schmidhuber2010}, the active-inference curiosity criterion \citep{FristonEtAl2017Curiosity}, and the information-theoretic bounded-rationality criterion \citep{OrtegaBraun2011,OrtegaBraun2013}. These traditions are not approximations of each other; they coincide on the action rule they prescribe. This is what makes the near-miss proved below substantive: Theorem~\ref{thm:afe-near-miss} is a failure of generic information-seeking itself, across a family of traditions that agree on what generic information-seeking means.
\end{remark}

\subsection{The algorithmic free energy principle}

\begin{definition}[Algorithmic free energy principle]
\label{def:afe-principle}
A learner obeys the \emph{algorithmic free energy principle} if
\begin{itemize}
\item[(i)] its prior $w$ is the universal interactive prior of Definition~\ref{def:universal-prior};
\item[(ii)] its posterior $q_t^\star(\cdot\mid h_t)$ is the Bayes/Gibbs posterior of Lemma~\ref{lem:variational};
\item[(iii)] its action $A_{t+1}$ is chosen to minimize the expected free energy $\Gcal_t(\cdot\mid h_t)$ with zero action cost and uniform preferences on a common finite feasible observation set, breaking ties by a fixed measurable rule.
\end{itemize}
\end{definition}

\subsection{Information decomposition relative to $[p^\star]$}

\begin{lemma}[Information decomposition]
\label{lem:info-decomp}
Fix a history $h_t$ and an action $a$. Let $c^\star=[p^\star]$ be the true interventional semantic class, let $\alpha_t:=q_t^\star(c^\star\mid h_t)\in(0,1)$, and let $Z:=\one\{P\in c^\star\}$ under $P\sim q_t^\star(\cdot\mid h_t)$. Then
\[
\I(P;O\mid h_t,a)
=
\I(Z;O\mid h_t,a)+(1-\alpha_t)\,\I(P;O\mid Z=0,h_t,a).
\]
\end{lemma}

\begin{proof}
Because $Z$ is a function of $P$, the chain rule gives
\[
\I(P;O\mid h_t,a)=\I(Z;O\mid h_t,a)+\I(P;O\mid Z,h_t,a).
\]
All programs in $c^\star$ induce the same next-step observation law, so $\I(P;O\mid Z=1,h_t,a)=0$. Taking the mixture over $Z$ at weights $\alpha_t,1-\alpha_t$ yields the claim.
\end{proof}

\begin{remark}
Lemma~\ref{lem:info-decomp} isolates the gap. Only the first term, information about membership in $[p^\star]$, is load-bearing for interventional semantic recovery. The second term can refine distinctions wholly inside $R_\pi(p^\star)\setminus[p^\star]$.
\end{remark}

\subsection{Near-miss: the AFE principle need not converge to $[p^\star]$}

Prior work on AIXI established that universal-prior learners can fail to identify the true environment in various senses \citep{LeikeHutter2015,Leike2016}, and that repair requires a specifically-chosen action rule such as Thompson sampling \citep{LeikeEtAl2016}. Theorem~\ref{thm:afe-near-miss} continues that line. It locates the failure of the same Bayes/Gibbs/expected-free-energy architecture, under a substituted finite-support prior, specifically at the level of the interventional semantic class $[p^\star]$, and shows that the expected-free-energy action rule is not a qualifying repair: the posterior on $[p^\star]$ can be frozen at its prior value for arbitrarily long horizons.

\begin{theorem}[Near-miss: the algorithmic free energy principle does not imply semantic convergence]
\label{thm:afe-near-miss}
Assume $\Ocal$ is infinite and $|\Acal|\ge 2$. For every $\alpha_0\in(0,1/2)$ and every horizon $T\ge 1$, there exist a computable interactive environment $p^\star\in\Pcal$ and a proper prior $w$ on $\Pcal$ supported on finitely many computable programs, with induced class prior $\bar w([p^\star])=\alpha_0$, such that the algorithmic free energy principle (Definition~\ref{def:afe-principle}) using $w$ in place of the universal interactive prior plays $a^{\mathrm{ref}}$ at every time $1,\dots,T$ on every sample path, and
\[
q_t^\star([p^\star]\mid H_t)=\alpha_0
\qquad\text{for every }t\in\{0,1,\dots,T\}.
\]
Since $T$ may be taken arbitrarily large and $\alpha_0\in(0,1/2)$ arbitrary, no rate of posterior convergence on $[p^\star]$ holds uniformly across algorithmic free energy learning problems.
\end{theorem}

The proof chooses $N+2$ distinct observations and two action modes: a ``separating'' action that discriminates $[p^\star]$ from its complement with at most $\log 2$ nats of mutual information, and a ``refining'' action that discriminates hypotheses within the complement of $[p^\star]$ with exactly $(1-\alpha_0)\log N$ nats. For $N$ large, expected free energy always prefers the refining action.

\begin{proof}
\emph{Construction.} Fix an integer $N\ge 3$ to be chosen below and the horizon $T\ge 1$. Choose distinct observations $R_1,\dots,R_N,S_0,S_1\in\Ocal$ and distinct actions $a^{\mathrm{ref}},a^{\mathrm{sep}}\in\Acal$.

Let $\Sigma_T:=\{1,\dots,N\}^T$. For each $\sigma\in\Sigma_T$, fix a computable program $p_\sigma\in\Pcal$ with action-conditioned law
\[
\mu_{p_\sigma}(o\mid h_{t-1},a)
=
\begin{cases}
\one\{o=R_{\sigma_t}\}, & a=a^{\mathrm{ref}},\ 1\le t\le T,\\
\one\{o=R_1\},         & a=a^{\mathrm{ref}},\ t>T,\\
\one\{o=S_0\},         & a\neq a^{\mathrm{ref}},
\end{cases}
\]
where the time index $t$ is read off from the length of $h_{t-1}$. Fix a computable program $p^\star\in\Pcal$ with law
\[
\mu_{p^\star}(o\mid h_{t-1},a)
=
\begin{cases}
N^{-1}\one\{o\in\{R_1,\dots,R_N\}\}, & a=a^{\mathrm{ref}},\\
\one\{o=S_1\},                        & a\neq a^{\mathrm{ref}}.
\end{cases}
\]
All such programs lie in $\Pcal$ under the universal prefix machine.

Write $c^\star:=[p^\star]$ and $c_\sigma:=[p_\sigma]$ for $\sigma\in\Sigma_T$. The classes $c^\star$ and $c_\sigma$ are distinct because the two programs disagree on $a^{\mathrm{sep}}$; classes $c_\sigma,c_{\sigma'}$ are distinct when $\sigma\neq\sigma'$ because the programs disagree at the first $t$ where $\sigma_t\neq\sigma'_t$ under $a^{\mathrm{ref}}$.

\emph{Prior.} Define a proper prior $w$ on $\Pcal$ by
\[
w(p^\star):=\alpha_0,\qquad w(p_\sigma):=\frac{1-\alpha_0}{N^T}\text{ for each }\sigma\in\Sigma_T,
\]
and $w\equiv 0$ on all other programs. This is a probability measure supported on $1+N^T$ computable programs. The induced class prior satisfies $\bar w(c^\star)=\alpha_0$ and $\bar w(c_\sigma)=(1-\alpha_0)/N^T$.

\emph{Learner.} Fix the algorithmic free energy principle of Definition~\ref{def:afe-principle} with the prior $w$ in place of the universal interactive prior, zero action cost, uniform preferences on the finite set $\{R_1,\dots,R_N,S_0,S_1\}$, and ties in action selection broken in favor of $a^{\mathrm{ref}}$.

\emph{Posterior on $c^\star$ under $a^{\mathrm{ref}}$-only transcripts.} Suppose the learner has played $a^{\mathrm{ref}}$ at times $1,\dots,t$ with observations $o_1,\dots,o_t\in\{R_1,\dots,R_N\}$ and $t\le T$. Encode $R_j\mapsto j$ to obtain $u\in\{1,\dots,N\}^t$. The number of $\sigma\in\Sigma_T$ matching $u$ on the first $t$ coordinates is exactly $N^{T-t}$, each with prior $(1-\alpha_0)/N^T$ and likelihood $1$ on the realized prefix. The program $p^\star$ has likelihood $N^{-t}$. Hence
\begin{equation}
\label{eq:nearmiss-alpha}
q_t^\star(c^\star\mid H_t)=\frac{\alpha_0\cdot N^{-t}}{\alpha_0\cdot N^{-t}+N^{T-t}\cdot (1-\alpha_0)/N^T}=\frac{\alpha_0}{\alpha_0+(1-\alpha_0)}=\alpha_0.
\end{equation}
The posterior on $c^\star$ equals its prior, on every sample path, as long as $t\le T$ and only $a^{\mathrm{ref}}$ has been played.

\emph{Conditional predictive among surviving $c_\sigma$'s.} The set of surviving $c_\sigma$'s at time $t$ is $\{\sigma\in\Sigma_T:\sigma_{1:t}=u\}$, of cardinality $N^{T-t}$. Because surviving classes have equal prior and equal likelihood, the posterior is uniform over this set. For $0\le t<T$, the conditional probability of $\sigma_{t+1}=j$ given the surviving set is
\[
\beta_t(j):=\frac{|\{\sigma\in\Sigma_T:\sigma_{1:t}=u,\,\sigma_{t+1}=j\}|}{N^{T-t}}=\frac{N^{T-t-1}}{N^{T-t}}=\frac{1}{N},
\qquad j\in\{1,\dots,N\}.
\]
So $\beta_t$ is \emph{exactly} uniform on $\{1,\dots,N\}$ for every $t\in\{0,1,\dots,T-1\}$ and every realized prefix $u$.

\emph{Mutual information at $a^{\mathrm{ref}}$.} Under $P=p^\star$ the next observation is uniform on $\{R_1,\dots,R_N\}$, so $\Hh(O\mid P=p^\star,h_t,a^{\mathrm{ref}})=\log N$. Under $P=p_\sigma$ (for any surviving $\sigma$) the next observation is $R_{\sigma_{t+1}}$ deterministically, so $\Hh(O\mid P=p_\sigma,h_t,a^{\mathrm{ref}})=0$. By \eqref{eq:nearmiss-alpha},
\[
\Hh(O\mid P,h_t,a^{\mathrm{ref}})=\alpha_0\log N+(1-\alpha_0)\cdot 0=\alpha_0\log N.
\]
The marginal predictive on $\{R_1,\dots,R_N\}$ is
\[
\lambda_t(j)=\alpha_0\cdot\frac{1}{N}+(1-\alpha_0)\cdot\beta_t(j)=\alpha_0\cdot\frac{1}{N}+(1-\alpha_0)\cdot\frac{1}{N}=\frac{1}{N},
\]
so $\Hh(O\mid h_t,a^{\mathrm{ref}})=\log N$. Therefore
\begin{equation}
\label{eq:mi-ref}
\I(P;O\mid h_t,a^{\mathrm{ref}})=\log N-\alpha_0\log N=(1-\alpha_0)\log N.
\end{equation}

\emph{Mutual information at $a^{\mathrm{sep}}$.} Under $a^{\mathrm{sep}}$, the observation is $S_1$ exactly when $P\in c^\star$ and $S_0$ otherwise. With $Z:=\one\{P\in c^\star\}$ under the posterior $q_t^\star$, the observation is a deterministic function of $Z$, so
\begin{equation}
\label{eq:mi-sep}
\I(P;O\mid h_t,a^{\mathrm{sep}})=\I(Z;O\mid h_t,a^{\mathrm{sep}})=\Hh(Z)=H_2(\alpha_0)\le\log 2,
\end{equation}
where $H_2(x):=-x\log x-(1-x)\log(1-x)$.

\emph{AFE plays $a^{\mathrm{ref}}$ through time $T$.} With zero action cost and uniform preferences on $\{R_1,\dots,R_N,S_0,S_1\}$, Lemma~\ref{lem:risk-ig} reduces expected-free-energy minimization to maximization of the mutual information $\I(P;O\mid h_t,a)$. Choose $N$ large enough that
\begin{equation}
\label{eq:nearmiss-N}
N>2^{1/(1-\alpha_0)},\qquad\text{equivalently }(1-\alpha_0)\log N>\log 2.
\end{equation}
Combining \eqref{eq:mi-ref}, \eqref{eq:mi-sep}, and \eqref{eq:nearmiss-N}, at every time $t\in\{0,\dots,T-1\}$ on which the learner has played $a^{\mathrm{ref}}$ exclusively,
\[
\I(P;O\mid h_t,a^{\mathrm{ref}})=(1-\alpha_0)\log N>\log 2\ge H_2(\alpha_0)=\I(P;O\mid h_t,a^{\mathrm{sep}}).
\]
The AFE rule selects $a^{\mathrm{ref}}$ (strictly; no tie-breaking is needed). Induction on $t$ shows that on every sample path, the learner plays $a^{\mathrm{ref}}$ at every time $1,\dots,T$.

\emph{Conclusion.} By \eqref{eq:nearmiss-alpha}, on every sample path,
\[
q_t^\star([p^\star]\mid H_t)=\alpha_0<1
\qquad\text{for every }t\in\{0,1,\dots,T\}.
\]
Since $T$ and $\alpha_0\in(0,1/2)$ were arbitrary, no rate of posterior convergence on $[p^\star]$ is uniform across AFE principle learning problems.
\end{proof}

\begin{theorem}[Observer-promotion failure of the AFE principle]
\label{thm:observer-promotion-failure}
Let $\pi^\sharp$ be the fixed policy that plays $a^{\mathrm{ref}}$ at every history. In the near-miss construction of Theorem~\ref{thm:afe-near-miss}, $\pi^\sharp$ agrees with the AFE principle's realized actions on every sample path through horizon $T$, and
\[
\omega_{\pi^\sharp}<\omega_{\mathrm{behav}}
\]
strictly in the refinement preorder.
\end{theorem}

\begin{proof}
Define $p'\in\Pcal$ by
\[
\mu_{p'}(o\mid h,a)=
\begin{cases}
\mu_{p^\star}(o\mid h,a^{\mathrm{ref}}), & a=a^{\mathrm{ref}},\\
\one\{o=S_0\}, & a\neq a^{\mathrm{ref}}.
\end{cases}
\]
Since $p^\star$'s $a^{\mathrm{ref}}$-law is computable, overriding its $a^{\mathrm{sep}}$-conditional with a fixed point-mass yields another program in $\Pcal$. Under $\pi^\sharp$ only $a^{\mathrm{ref}}$-pairs are reachable, and $\mu_{p'}$ agrees with $\mu_{p^\star}$ there, so $\Prob_{p',\pi^\sharp}=\Prob_{p^\star,\pi^\sharp}$ by the history-factorization identity; hence $p'\in R_{\pi^\sharp}(p^\star)$. The pair $(\varepsilon,a^{\mathrm{sep}})$ is jointly reachable, and $\mu_{p'}(\cdot\mid\varepsilon,a^{\mathrm{sep}})=\delta_{S_0}\neq\delta_{S_1}=\mu_{p^\star}(\cdot\mid\varepsilon,a^{\mathrm{sep}})$, so the kernel-equality convention of Section~\ref{sec:setup} gives $\mu_{p'}\neq\mu_{p^\star}$, i.e., $p'\notin[p^\star]$. Hence $R_{\pi^\sharp}(p^\star)\supsetneq[p^\star]$, which by Definition~\ref{def:observer} is $\omega_{\pi^\sharp}<\omega_{\mathrm{behav}}$.
\end{proof}

\begin{corollary}[Universal-prior case]
\label{cor:observer-promotion-universal}
The strict refinement gap $\omega_{\pi^\sharp}<\omega_{\mathrm{behav}}$ of Theorem~\ref{thm:observer-promotion-failure} is prior-independent. Replacing the near-miss's finite-support prior by the universal interactive prior does not alter this gap for the fixed policy $\pi^\sharp$.
\end{corollary}

\begin{proof}
The witness $p'$ and the policy $\pi^\sharp$ in Theorem~\ref{thm:observer-promotion-failure} are defined without reference to the prior. The same argument therefore yields the same strict refinement gap for $\pi^\sharp$ under the universal interactive prior.
\end{proof}

\begin{remark}
\label{rem:near-miss}
Theorem~\ref{thm:afe-near-miss} strengthens the near-miss from a one-step counterexample to an arbitrary-horizon convergence-failure family. Theorem~\ref{thm:observer-promotion-failure} reads this as a structural failure of action-side observer promotion: the AFE principle's realized observer $\omega_{\pi^\sharp}$ sits strictly below $\omega_{\mathrm{behav}}$, and the frozen-posterior conclusion is its Bayesian shadow. The belief rule is exact and the prior assigns positive mass to the true class; the failure is entirely on the action side, and a stricter action primitive --- the one Section~\ref{sec:main} supplies --- is required.
\end{remark}

\begin{remark}
\label{rem:near-miss-universal}
Theorem~\ref{thm:afe-near-miss} is proved for a substituted finite-support prior. Extending its exact frozen-posterior conclusion to the literal universal interactive prior would require a separate robustness argument controlling how the universal prior's additional support perturbs both the posterior values and the expected-free-energy action ranking. Corollary~\ref{cor:observer-promotion-universal} therefore isolates the prior-independent part of the failure: the strict refinement gap for the fixed policy $\pi^\sharp$ persists under the universal prior.
\end{remark}

\section{Class-Against-Complement Action}
\label{sec:semantic-action}

The near-miss isolates the comparison that matters. Once the target is a class $c$, the relevant rival is not any single alternative program. It is all posterior mass outside $c$ pooled together. An action that fails to separate $c$ from that aggregate rival may be highly informative about the latent program and still supply no evidence for convergence to $c$.

For a class $c\in\Ccal$, write $\mu_c$ for the common law induced by any program in $c$.

\begin{definition}[Class-complement predictive law]
\label{def:class-complement}
Fix a history $h_t$, a class $c\in\Ccal$, and an action $a$. If $0<q_t^\star(c\mid h_t)<1$, define
\[
\Qcal_t^{-c}(o\mid h_t,a)
:=
\frac{\sum_{d\neq c} q_t^\star(d\mid h_t)\mu_d(o\mid h_t,a)}
{1-q_t^\star(c\mid h_t)}.
\]
\end{definition}

\begin{definition}[Semantic gain]
\label{def:semantic-gain}
The \emph{semantic gain} of action $a$ for class $c$ at $h_t$ is
\[
S_t(c,a;h_t)
:=
\begin{cases}
\KL\!\bigl(\mu_c(\cdot\mid h_t,a)\,\|\,\Qcal_t^{-c}(\cdot\mid h_t,a)\bigr), & 0<q_t^\star(c\mid h_t)<1,\\
0, & q_t^\star(c\mid h_t)\in\{0,1\}.
\end{cases}
\]
\end{definition}

\begin{lemma}[Posterior-odds identity]
\label{lem:odds-identity}
Fix $c\in\Ccal$ with $0<q_t^\star(c\mid h_t)<1$ and an action $a$. Let the latent truth lie in $c$. Then
\[
S_t(c,a;h_t)
=
\E_{o\sim\mu_c(\cdot\mid h_t,a)}\!\left[
\log\frac{q_{t+1}^\star(c\mid h_t,a,o)}{1-q_{t+1}^\star(c\mid h_t,a,o)}
-\log\frac{q_t^\star(c\mid h_t)}{1-q_t^\star(c\mid h_t)}
\right].
\]
\end{lemma}

\begin{proof}
Let $\alpha_t:=q_t^\star(c\mid h_t)$. Class-level Bayes gives
\[
\frac{q_{t+1}^\star(c\mid h_t,a,o)}{1-q_{t+1}^\star(c\mid h_t,a,o)}
=
\frac{\alpha_t}{1-\alpha_t}\cdot\frac{\mu_c(o\mid h_t,a)}{\Qcal_t^{-c}(o\mid h_t,a)}.
\]
Taking logarithms and expectation under $\mu_c(\cdot\mid h_t,a)$ yields the claim.
\end{proof}

\begin{definition}[Semantic separation]
\label{def:semantic-separation}
The \emph{semantic separation} of action $a$ for class $c$ at $h_t$ is
\[
B_t(c,a;h_t)
:=
\begin{cases}
-\log\sum_{o\in\Ocal}\sqrt{\mu_c(o\mid h_t,a)\Qcal_t^{-c}(o\mid h_t,a)}, & 0<q_t^\star(c\mid h_t)<1,\\
0, & q_t^\star(c\mid h_t)\in\{0,1\}.
\end{cases}
\]
\end{definition}

\begin{lemma}[Zero criterion]
\label{lem:zero-criterion}
For $c$ and $a$ with $0<q_t^\star(c\mid h_t)<1$,
\[
S_t(c,a;h_t)=0
\iff
B_t(c,a;h_t)=0
\iff
\mu_c(\cdot\mid h_t,a)=\Qcal_t^{-c}(\cdot\mid h_t,a).
\]
\end{lemma}

\begin{proof}
KL vanishes exactly at coincidence of laws. The Hellinger affinity equals one exactly at coincidence. The zero set is preserved under $-\log$.
\end{proof}

\begin{definition}[Semantic action rule]
\label{def:semantic-rule}
Let $\bar w(c):=\sum_{p\in c}w(p)$ be the induced prior on classes. Fix $\epsilon\in(0,1]$. For each class $c$ and history $h_t$, choose a measurable selector $a_t^c(h_t)$ with
\[
\tfrac{1}{2}\min\!\Bigl\{1,\sup_{a\in\Acal}B_t(c,a;h_t)\Bigr\}
\le
B_t(c,a_t^c(h_t);h_t)
\]
when $0<q_t^\star(c\mid h_t)<1$, and arbitrary otherwise. Define the class-targeting distribution
\[
\nu_t(c\mid h_t):=(1-\epsilon)q_t^\star(c\mid h_t)+\epsilon\bar w(c).
\]
The \emph{semantic action rule} samples $C_t\sim\nu_t(\cdot\mid h_t)$ and plays $a_{t+1}=a_t^{C_t}(h_t)$.
\end{definition}

Such a selector exists because $\Acal$ is countable, so one may fix an enumeration and choose the first action satisfying the displayed inequality.

\begin{definition}[Promotion-supporting action rule]
\label{def:promotion-supporting}
An action rule is \emph{promotion-supporting} if there exists $\delta>0$ such that for every class $c\in\Ccal$ and every history $h_t$ with $0<q_t^\star(c\mid h_t)<1$,
\[
\Prob\bigl(A_{t+1}=a_t^c(h_t)\mid H_t=h_t\bigr)\ge \delta\bar w(c),
\]
where $a_t^c(h_t)$ is the $c$-separating selector of Definition~\ref{def:semantic-rule}.
\end{definition}

\begin{proposition}[Semantic rule is promotion-supporting]
\label{prop:semantic-is-promotion-supporting}
The semantic action rule of Definition~\ref{def:semantic-rule} is promotion-supporting with $\delta=\epsilon$. In particular, the true class $c^\star=[p^\star]$ satisfies
\[
\Prob(C_t=c^\star\mid H_t)\ge\epsilon\bar w(c^\star)>0
\]
almost surely for every $t$.
\end{proposition}

\begin{proof}
Fix $h_t$ and $c\in\Ccal$ with $0<q_t^\star(c\mid h_t)<1$. Conditioning on $C_t=c$,
\[
\Prob(A_{t+1}=a_t^c(h_t)\mid H_t=h_t)\ge\Prob(C_t=c\mid H_t=h_t)=\nu_t(c\mid h_t)\ge\epsilon\bar w(c).
\]
The universal prior assigns positive mass to every computable program, so $\bar w(c)>0$ for every $c\in\Ccal$.
\end{proof}

\section{Sufficient Conditions for Semantic Convergence}
\label{sec:main}

Section~\ref{sec:near-miss} located the algorithmic free energy principle's failure on the action side: in the near-miss construction, its realized policy induces an observer strictly coarser than $\omega_{\mathrm{behav}}$ (Theorem~\ref{thm:observer-promotion-failure}). This section supplies a class of action rules whose Bayesian effect is the opposite --- the posterior concentrates on $\omega_{\mathrm{behav}}$'s indistinguishability class $[p^\star]$, the observer-promotion success symmetric to the AFE near-miss.

\begin{definition}[Semantic separation condition]
\label{def:sep-condition}
The interactive setting satisfies the \emph{semantic separation condition} if for every class $c\in\Ccal$ there exists $\eta(c)>0$ such that whenever $0<q_t^\star(c\mid h_t)<1$,
\[
\sup_{a\in\Acal}B_t(c,a;h_t)\ge \eta(c).
\]
\end{definition}

\begin{lemma}[Conditional Borel--Cantelli]
\label{lem:conditional-bc}
Let $(\mathcal{F}_t)_{t\ge 0}$ be a filtration and $E_{t+1}\in\mathcal{F}_{t+1}$ for every $t\ge 0$. If $\sum_t\Prob(E_{t+1}\mid\mathcal{F}_t)=\infty$ almost surely, then $E_{t+1}$ i.o.\ almost surely.
\end{lemma}

\begin{proof}
See \citet[Chapter~12]{Williams1991}.
\end{proof}

\begin{lemma}[Cumulative separation drives posterior odds to zero]
\label{lem:contraction}
Fix $c^\star=[p^\star]$. Under the universal prior and Bayes/Gibbs posterior, if
\[
\sum_{t=0}^{\infty}B_t(c^\star,A_{t+1};H_t)=\infty
\qquad
\text{almost surely,}
\]
then $q_t^\star(c^\star\mid H_t)\to 1$ almost surely.
\end{lemma}

\begin{proof}
Let $\mathcal{G}_t:=\sigma(H_t,A_{t+1})$. Define the posterior odds against $c^\star$ by
\[
R_t^-:=\frac{1-q_t^\star(c^\star\mid H_t)}{q_t^\star(c^\star\mid H_t)}
=\frac{\sum_{p\notin c^\star}w(p)\mu_p(o_{1:t}\givena a_{1:t})}
      {\bar w(c^\star)\mu_{p^\star}(o_{1:t}\givena a_{1:t})}.
\]
Whenever $R_t^->0$, Definition~\ref{def:class-complement} gives
\[
R_{t+1}^-=R_t^-\cdot\frac{\Qcal_t^{-c^\star}(o_{t+1}\mid H_t,A_{t+1})}{\mu_{p^\star}(o_{t+1}\mid H_t,A_{t+1})}.
\]
Taking square roots and conditional expectation under the true law,
\[
\E\!\left[\sqrt{R_{t+1}^-}\mid\mathcal{G}_t\right]
=
\sqrt{R_t^-}\sum_{o\in\Ocal}\sqrt{\mu_{p^\star}(o\mid H_t,A_{t+1})\Qcal_t^{-c^\star}(o\mid H_t,A_{t+1})}
=
\sqrt{R_t^-}\,e^{-B_t(c^\star,A_{t+1};H_t)}.
\]
Let $S_t:=\sum_{i=0}^{t-1}B_i(c^\star,A_{i+1};H_i)$ and $M_t:=e^{S_t}\sqrt{R_t^-}$. Then $(M_t)$ is a nonnegative $\mathcal{G}_t$-martingale, hence converges almost surely. On $\{S_t\to\infty\}$, $\sqrt{R_t^-}=M_te^{-S_t}\to 0$, so $R_t^-\to 0$ and $q_t^\star(c^\star\mid H_t)\to 1$.
\end{proof}

\begin{theorem}[Sufficient conditions for semantic convergence]
\label{thm:semantic-convergence}
Assume:
\begin{itemize}
\item[(C1)] the prior is the universal interactive prior (Definition~\ref{def:universal-prior});
\item[(C2)] belief is the Bayes/Gibbs posterior (Lemma~\ref{lem:variational});
\item[(C3)] the semantic separation condition (Definition~\ref{def:sep-condition}) holds;
\item[(C4)] action follows the semantic action rule (Definition~\ref{def:semantic-rule}).
\end{itemize}
Then for every $p^\star\in\Pcal$,
\[
q_t^\star([p^\star]\mid H_t)\longrightarrow 1
\qquad
\text{almost surely.}
\]
\end{theorem}

\begin{proof}
Let $\mathcal{F}_t:=\sigma(H_t,C_0,\dots,C_{t-1})$. The event $E_t:=\{q_t^\star(c^\star\mid H_t)<1\}$ is $\mathcal{F}_t$-measurable, and $T_{t+1}:=E_t\cap\{C_t=c^\star\}$ is the event that the semantic action rule targets $c^\star$ while uncertainty remains. By Proposition~\ref{prop:semantic-is-promotion-supporting},
\[
\Prob(T_{t+1}\mid\mathcal{F}_t)
=\one_{E_t}\nu_t(c^\star\mid H_t)\ge\epsilon\bar w(c^\star)\one_{E_t}.
\]
Set $\hat\eta:=\tfrac{1}{2}\min\{1,\eta(c^\star)\}>0$. On $E_t$ i.o., $\sum_t\Prob(T_{t+1}\mid\mathcal{F}_t)=\infty$, so Lemma~\ref{lem:conditional-bc} gives $T_{t+1}$ i.o.\ almost surely on $E_t$ i.o. Whenever $T_{t+1}$ occurs, the selector satisfies $B_t(c^\star,A_{t+1};H_t)\ge\hat\eta$ by the semantic separation condition. Hence on $E_t$ i.o.,
\[
\sum_{t=0}^{\infty}B_t(c^\star,A_{t+1};H_t)=\infty
\qquad
\text{almost surely,}
\]
and Lemma~\ref{lem:contraction} gives $q_t^\star(c^\star\mid H_t)\to 1$ on that event. On $E_t^c$ eventually, the same conclusion is immediate.
\end{proof}

\begin{remark}
\label{rem:promotion-support}
The proof invokes the semantic action rule only through Proposition~\ref{prop:semantic-is-promotion-supporting}. Any action rule satisfying Definition~\ref{def:promotion-supporting} and (C1)--(C3) obeys the same conclusion. Theorem~\ref{thm:semantic-convergence} thus identifies the class-against-complement rule as one instance of a broader family of observer-promoting action rules, with the near-miss construction of Theorem~\ref{thm:afe-near-miss} furnishing a failure of promotion support on the AFE side.
\end{remark}

\begin{corollary}[Finite-class deterministic specialization]
\label{cor:finite-maximin}
Under (C1), (C2) and a finite semantic class set $\Ccal$ with $\eta>0$ such that whenever $q_t^\star(c^\star\mid H_t)<1$, some action $a$ satisfies $B_t(c,a;H_t)\ge\eta$ for every unresolved class $c$, define
\[
U_t(a;H_t):=
\begin{cases}
\min_{c:\,0<q_t^\star(c\mid H_t)<1}B_t(c,a;H_t), & q_t^\star(c^\star\mid H_t)<1,\\
0, & q_t^\star(c^\star\mid H_t)=1,
\end{cases}
\]
and whenever $q_t^\star(c^\star\mid H_t)<1$ let $A_{t+1}$ satisfy $U_t(A_{t+1};H_t)\ge\tfrac{1}{2}\min\{1,\sup_a U_t(a;H_t)\}$, choosing $A_{t+1}$ arbitrarily when $q_t^\star(c^\star\mid H_t)=1$. Then $q_t^\star(c^\star\mid H_t)\to 1$ almost surely.
\end{corollary}

\begin{proof}
Whenever the posterior is not concentrated, $\sup_a U_t(a;H_t)\ge\eta$, so $B_t(c^\star,A_{t+1};H_t)\ge\tfrac{1}{2}\min\{1,\eta\}$ whenever any false posterior mass remains. Cumulative separation diverges and Lemma~\ref{lem:contraction} applies.
\end{proof}

\begin{corollary}[Support on separating interventions is necessary]
\label{cor:support-necessary}
If a policy class permits zero probability on a semantically separating action at a reachable history, then no uniform theorem of concentration on $[p^\star]$ can hold over $\Pcal$ for that policy class under any Bayesian scheme whose prior assigns positive mass to every computable program.
\end{corollary}

\begin{proof}
Theorem~\ref{thm:policy-gap} produces $p'\in R_\pi(p^\star)\setminus[p^\star]$. For every horizon $T$ and every history $h_T$ with $\Prob_{p^\star,\pi}(H_T=h_T)>0$, every program in $[p^\star]$ has the realized likelihood of $p^\star$ and every program in $[p']$ has the realized likelihood of $p'$; because $p'\in R_\pi(p^\star)$, these class likelihoods are equal under $\pi$. Hence
\[
\frac{q_T([p']\mid h_T)}{q_T([p^\star]\mid h_T)}
=
\frac{\bar w([p'])}{\bar w([p^\star])}
>0.
\]
Therefore
\[
q_T([p^\star]\mid h_T)\le \frac{1}{1+\bar w([p'])/\bar w([p^\star])}<1
\]
at every horizon on the true-history support, so concentration on $[p^\star]$ fails.
\end{proof}

\begin{corollary}[Observer-promotion contrast]
\label{cor:promotion-contrast}
In the near-miss construction of Theorem~\ref{thm:afe-near-miss}, replace the prior by one satisfying $w(p^\star)=\alpha_0$, $w(p')\in(0,1-\alpha_0)$ for the witness $p'$ of Theorem~\ref{thm:observer-promotion-failure}, and $w(p_\sigma)=(1-\alpha_0-w(p'))/N^T$ for each $\sigma\in\Sigma_T$, and choose $N>2^{1/(1-\alpha_0-w(p'))}$ in the construction. Then the posterior ratio $q_t^\star(p'\mid H_t)/q_t^\star(p^\star\mid H_t)$ is constant and equal to $w(p')/\alpha_0$ on every sample path for $t\le T$ under the AFE principle, while $q_t^\star(p'\mid H_t)\to 0$ almost surely under any promotion-supporting action rule satisfying (C1), (C2), and (C3).
\end{corollary}

\begin{proof}
Through horizon $T$ the AFE rule plays $\pi^\sharp$: the strengthened choice $N>2^{1/(1-\alpha_0-w(p'))}$ maintains $(1-\alpha_0-w(p'))\log N>\log 2\ge H_2(\alpha_0)$, so expected-free-energy's mutual information at $a^{\mathrm{ref}}$ strictly exceeds that at $a^{\mathrm{sep}}$ as in Theorem~\ref{thm:afe-near-miss}'s proof applied to the enlarged prior. Since $\mu_{p'}=\mu_{p^\star}$ at every reachable pair under $\pi^\sharp$, realized-history likelihoods coincide and the posterior ratio equals the prior ratio. For the promotion side, Theorem~\ref{thm:semantic-convergence} and Remark~\ref{rem:promotion-support} give $q_t^\star([p^\star]\mid H_t)\to 1$ almost surely under any promotion-supporting action rule satisfying (C1), (C2), and (C3). Because $p'\notin[p^\star]$, one has $q_t^\star(p'\mid H_t)\le 1-q_t^\star([p^\star]\mid H_t)\to 0$ almost surely.
\end{proof}

\begin{remark}
\label{rem:promotion-contrast-interpretation}
Corollary~\ref{cor:promotion-contrast} reads the knowability gap $[p^\star]\subsetneq R_{\pi^\sharp}(p^\star)$ as an action-rule test. The AFE rule respects the gap and cannot expel $p'$ from the posterior; a promotion-supporting rule probes the gap via $a^{\mathrm{sep}}$ and resolves it. Theorem~\ref{thm:policy-gap}'s set inclusion is exactly the structural object an observer-promoting action rule must close.
\end{remark}

\section{Implications}
\label{sec:implications}

The four-set hierarchy clarifies where existing learning theories stand. Each theory targets one of the four observers of Definition~\ref{def:observer}; the question for any given theory is whether its learning rule promotes the realized observer to the observer whose indistinguishability class the theory tries to identify.

Realized-history fit aims at $R_{h_t}(p^\star)$. Next-token training objectives under any fixed data distribution are instances: the learner scores well by assigning high likelihood to realized prefixes. Lemma~\ref{lem:fit-gap} is the boundary: such a learner's posterior can be on a compatibility set whose next-queried continuation is still undetermined. Predictive merging (Lemma~\ref{lem:merging}) promotes the posterior to the one-step-predictive interior of $R_{h_t}(p^\star)$ along the realized interaction.

Policy-relative prediction aims at $R_\pi(p^\star)$. Off-policy reinforcement learning and partially-observed Markov decision processes without active interventions face the boundary of Theorem~\ref{thm:policy-gap}: the posterior can concentrate on $R_\pi(p^\star)$ while leaving $[p^\star]$ unresolved.

Interventional semantic recovery aims at $[p^\star]$. It is the target of the main theorem and the theorem-level target of explanatory learning. The algorithmic free energy principle of Definition~\ref{def:afe-principle} is an exact universal-prior Bayesian learner with a generic epistemic action rule. Theorem~\ref{thm:afe-near-miss} shows that the same Bayes/Gibbs/expected-free-energy architecture, under a substituted finite-support prior, does not suffice to reach this target. Section~\ref{sec:semantic-action} identifies the stricter action primitive that does: increase the posterior odds of a semantic class against its complement. Concrete systems built around program-level structure --- Bayesian program learning and program induction \citep{LakeEtAl2015,EllisEtAl2021DreamCoder}, abstraction-and-reasoning benchmarks \citep{Chollet2024ARC}, and search over program spaces guided by large language models \citep{RomeraParedesEtAl2024FunSearch} --- can be read as instances of class-targeted learning in this sense, their semantic classes determined by the structure of the program space they search.

Canonical syntactic recovery aims at $\{p^\star\}$. Theorem~\ref{thm:factor-through-quotient} rules it out unless the designated representative is already constant on $[p^\star]$.

\section{Discussion}
\label{sec:discussion}

The boxed problem has two logically distinct parts, and the paper resolves them at different levels. The belief rule is fully characterized: the Bayes/Gibbs posterior under the universal prior is the unique minimizer of the algorithmic free energy (Lemma~\ref{lem:variational}), and within a natural convex divergence class the KL-based functional is the only one with that property (Lemma~\ref{lem:kl-necessity}). The action rule is more revealing. The theorem-level near-miss family shows that the same Bayes/Gibbs/expected-free-energy architecture, under a substituted finite-support prior, does not in general reach the interventional semantic class (Theorem~\ref{thm:afe-near-miss}). The action objective that supports Theorem~\ref{thm:semantic-convergence} must be class-targeted.

The observer frame places these action-rule results alongside three traditions. Pearl's do-calculus identifies when interventional distributions are determined by observational ones via graphical structure \citep{Pearl2009}; Corollary~\ref{cor:promotion-contrast} is its algorithmic analogue, naming the action-rule condition under which active play determines the interventional semantic class. \citet{Chernoff1959}'s asymptotically optimal sequential-testing rule selects, at each step, the design maximizing the separation of the most-likely hypothesis from the nearest alternative. The class-against-complement rule is its countable-hypothesis, universal-prior extension: promotion support (Definition~\ref{def:promotion-supporting}) plays the role of Chernoff's selection-probability condition, and the aggregate-complement law $\Qcal_t^{-c}$ replaces Chernoff's pairwise closest alternative when the rival class is structurally heterogeneous. \citet{Blackwell1953}'s comparison of experiments orders statistical designs by a sufficiency-type relation --- the parametric ancestor of Definition~\ref{def:observer}'s refinement preorder --- so Theorem~\ref{thm:semantic-convergence} reads as observer promotion in Blackwell's order along the learned trajectory.

The observer frame also sharpens three philosophical theses about underdetermination. \citet{Quine1960}'s claim that empirical evidence leaves logically incompatible theories equally supported is made formal by Theorem~\ref{thm:policy-gap}: whenever a fixed policy leaves some reachable history-action pair unqueried, there are concrete programs in $R_\pi(p^\star)\setminus[p^\star]$ that no amount of passive observation expels. \citet{vanFraassen1980}'s constructive empiricism holds that science aims at empirical adequacy rather than truth; Theorem~\ref{thm:factor-through-quotient} identifies the interactive analogue of empirical adequacy as $[p^\star]$, the finest target any history-based estimator can reliably recover. \citet{Putnam1980} questioned whether use could fix reference at all; Theorem~\ref{thm:semantic-convergence} shows that class-against-complement action fixes reference up to $[p^\star]$, though no farther. The observer formalism separates two kinds of underdetermination: the passive-active gap $[p^\star]\subsetneq R_\pi(p^\star)$, which active intervention closes via observer promotion, and the residue inside $[p^\star]$, which Theorem~\ref{thm:factor-through-quotient} shows no history-based observer can resolve.

Three theoretical questions remain open.

First, when do generic information gain and semantic gain coincide? Theorem~\ref{thm:afe-near-miss} shows that coincidence can fail, already in the substituted-prior family studied here. It is natural to ask for structural conditions on the hypothesis class under which expected-free-energy maximization automatically maximizes semantic gain. The answer would identify the precise boundary between generic epistemic discovery and explanatory action.

Second, what quantitative rates can be proved? Theorem~\ref{thm:semantic-convergence} is asymptotic. A finite-horizon lower bound on $q_t^\star([p^\star]\mid h_t)$ in terms of cumulative semantic separation would sharpen both theory and practice. Sequential design results in statistics suggest such bounds should be possible in restricted regimes \citep{Chernoff1959,BoxHill1967}.

Third, what tractable approximations preserve the theorem-level target? The universal prior and exact Bayes/Gibbs posterior are incomputable. The relevant approximation question is not whether a heuristic reduces uncertainty in general, but whether it preserves class-targeted semantic contraction.

\section{Conclusion}
\label{sec:conclusion}

This paper organizes interactive learning around four nested sets of programs, the indistinguishability classes of four observers at increasing observational power. Each inclusion is strict in general. The belief problem is solved exactly by the universal prior and Bayes/Gibbs posterior, which minimize the algorithmic free energy. The generic action objective, expected free energy, is a near-miss: in the theorem's substituted-prior family, the corresponding learner can induce a realized observer strictly coarser than the behavioral one, and its posterior on $[p^\star]$ need not place mass one on the true interventional semantic class. Concentration on $[p^\star]$ requires action-side observer promotion, achieved by a class-against-complement action primitive. Under the universal prior, Bayes/Gibbs belief, a semantic separation condition, and the semantic action rule, the posterior mass on $[p^\star]$ converges almost surely to one.

\bibliographystyle{plainnat}
\bibliography{algorithmic_free_energy_principle}

\end{document}
